\documentclass[11pt,letterpaper]{article}
\usepackage[margin=0.95in]{geometry}
\usepackage{amsmath,amssymb,amsthm,booktabs,longtable,array,xcolor}
\usepackage{hyperref}
\hypersetup{colorlinks=true,linkcolor=blue!45!black,citecolor=blue!45!black,
 urlcolor=blue!45!black,pdftitle={A formal zeta(9) analogue of Fauzan's construction},
 pdfauthor={Independent exploration prepared with Codex and GPT-6 Astra}}
\numberwithin{equation}{section}
\setlength{\emergencystretch}{2em}
\setlength{\parskip}{0.25em}
\allowdisplaybreaks[2]
\newcommand{\pos}[1]{\left(#1\right)_+}
\title{A formal $\zeta(9)$ analogue\\of Fauzan's construction}
\author{A section-by-section working draft}
\date{14 November 2026}
\begin{document}
\maketitle
\begin{abstract}
We rewrite the algebraic construction in Aabir Fauzan's September 2026
manuscript for $\zeta(9)$, retaining its determinant size, numerator
$D_N^6$, and scalar normalization. The ninth-order functional has an
explicit positive weight; its determinant has the same degree, and the
leading real energy calculation is unchanged. The local arithmetic is
more expensive. The normalization calculation in this draft increases
the quadratic arithmetic bound by exactly $116581/57600$. Its sum with
the reused real bound is positive, so these estimates supply no decay
statement from which irrationality would follow. Sections 1--7 and
Appendices A--C follow the source's organization, with local derivations,
formal transplants, reused numerical inputs, and blocked conclusions
identified separately. This is an exploratory analogue, not a proof of
irrationality or a verification of the source paper.
\end{abstract}

\smallskip
\noindent\textbf{How to read this draft.}
\textbf{[Derived]} means that the displayed local calculation was worked
out for this draft; it does not certify all surrounding arguments.
\textbf{[Formal analogue]} identifies a transplanted step whose full
justification has not been supplied.
\textbf{[Conditional]} means that the conclusion uses an explicitly
stated estimate or imported certificate.
\textbf{[Blocked conclusion]} marks a conclusion for which the required
sign or hypothesis is missing.

\noindent References such as ``source (4.12)'' refer to Fauzan's equation
numbers. Other equation references refer to this draft. Numerical
certificates retained from the source are identified as inputs rather
than as computations rerun here.

\tableofcontents
\clearpage
\section{Introduction and the proposed analogue}
\label{s1-introduction}

The source manuscript \cite{F} constructs rational polynomials in a
variable $X$, uses local estimates to clear their coefficients, and uses
a positive integral to estimate their values at $X=\zeta(5)$. Here we
carry out the corresponding manipulations with $\zeta(9)$. The question
is what the individual formulas become, including the formulas at the
end that cease to give the desired conclusion.

We make one choice at the outset: retain the source's
$K=40n$, $N=3n$, $h=37n$, and numerator $D_N^6$. Replace the functional
and all harmonic sums, moments, and residues that depend on the zeta
order. A construction with $D_N^{10}$ would be a different experiment:
it changes the degree bookkeeping, local weights, normalization, and
real field together. It is not part of the main calculation below.

\subsection{What the final step would require}

The mechanism behind the source's main theorem can be stated without
assuming any of its estimates. Suppose integer polynomials $Q_K$ have
degree $h=\lambda K$, where $\lambda=37/40$, and satisfy
\begin{equation}
 0<Q_K(\zeta(9))\leq \exp(-\varepsilon K^2+o(K^2))
 \quad\text{for some }\varepsilon>0.
 \label{s1-target}
\end{equation}
If $\zeta(9)=a/b$ with integers $a,b$ and $b>0$, then
$b^hQ_K(a/b)$ is a positive integer. But its logarithm is at most
$\lambda K\log b-\varepsilon K^2+o(K^2)$, which tends to $-\infty$.
That is the contradiction one would seek. Positivity, integrality, and
quadratic decay have distinct jobs in this short final argument.

Our analogue writes $Q_{9,K,M}=m_{9,K,M}F_{9,K}$, with $F_{9,K}$ the
real-normalized determinant and $m_{9,K,M}>0$ the arithmetic multiplier.
For fixed $M$, the two proposed bounds are
\begin{equation}
 \log m_{9,K,M}\leq A_{9,M}K^2+o_M(K^2),\qquad
 \log F_{9,K}(\zeta(9))\leq UK^2+O(K\log K).
 \label{s1-two-bounds}
\end{equation}
The sign that matters is therefore $A_{9,M}+U$. An upper bound with a
positive coefficient does not imply that the determinant values grow;
it simply fails to prove the decay needed in \eqref{s1-target}.

\subsection{The changes to expect}

The following dictionary records the principal substitutions. The
entries in the right column are derived or qualified in the indicated
sections; the table is a guide to the calculations, not an additional
theorem.
\begin{center}
\small
\begin{tabular}{@{}>{\raggedright\arraybackslash}p{0.35\textwidth}>{\raggedright\arraybackslash}p{0.25\textwidth}>{\raggedright\arraybackslash}p{0.29\textwidth}@{}}
\toprule
Object & Source order five & Draft order nine\\
\midrule
Simple-pole coefficient of $X$ & $j^4$ & $j^8$\\
Polynomial functional & $L(P^{(3)})/4!$ & $L(P^{(7)})/8!$\\
Real weight series & $y^5\sum \ell^4e^{-2\pi\ell y}$ & $y^9\sum \ell^8e^{-2\pi\ell y}$\\
Large-prime distribution & $p^{-4}$ & $p^{-8}$\\
Ordinary-class exponent loss & $-4$ & $-8$\\
Zero-class net contribution & $+1$ & $+1$\\
First possible nonintegral moment, $p\geq11$ & $e\geq2p-3$ & $e\geq4p-5$\\
Outer unremoved-class cost & $7(\ell-2)$ & $11(\ell-2)$\\
Degree of the determinant & $h$ & $h$\\
Leading real bound $U$ & $-1.3669955$ & $-1.3669955$\\
\bottomrule
\end{tabular}
\end{center}

Two effects compete in the arithmetic. Higher harmonic sums increase
the pole denominators. Extra factors in the polynomial moments cancel
more Bernoulli denominators. Section~4 keeps both effects: in particular,
the outer polynomial correction disappears sooner, even though the
total denominator-clearing estimate becomes more costly.

The calculation in Section~5 and Appendix~B gives
\begin{equation}
 A_{9,M}=A_{5,M}+\frac{116581}{57600},\qquad
 \frac{116581}{57600}=2.0239756944\ldots .
 \label{s1-cost-shift}
\end{equation}
Here $A_{5,M}$ denotes the source's particular arithmetic bound. It is
not a claimed minimum among all normalizations. Combining the draft's
bound with $U$ leaves approximately $2.00657$ at $M=200$, and approaches
$1.97449$ as $M\to\infty$. Thus this version of the estimates has already
lost its final negative sign. Section~7 nevertheless writes out the
completion, and Appendix~C carries the perturbation and height algebra
as far as its hypotheses permit.

There is no additional restriction $9\equiv1\pmod4$ in the positive
weight calculation used here. For every odd $s\geq3$, the corresponding
weight has the positive series
\[
 \frac{2(2\pi)^{s-1}}{(s-1)!}y^s
       \sum_{\ell\geq1}\ell^{s-1}e^{-2\pi\ell y}.
\]
Both the seventh and ninth orders fit that identity. The relevant
obstruction in this draft is the numerical arithmetic cost, not a
change of sign in that series. This statement concerns the present
functional; it makes no assertion about separate modular transformation
formulas in which congruence modulo four can matter.

\subsection{Organization and limits of the exercise}

Sections~2 and 3 construct the functional and translate its local
identities. Section~4 supplies the proposed determinant exponents.
Section~5 converts them to a multiplier and a prime-sum cost.
Section~6 preserves the real energy calculation while changing the
weight and its subquadratic errors. Section~7 displays the failed sign
comparison. Appendix~A shows how the source's comparison measure is
reused; Appendix~B gives the new rational arithmetic totals;
Appendix~C develops perturbation, rational-approximation, and height
formulas, retaining a hypothetical decay parameter only where needed.

The inner completion step remains marked as a formal transplant. The
source's real interval certificate is reused, not reverified. These
qualifications do not prevent one from examining the substitutions,
but they do prevent treating the draft as a checked theorem. The
positive combined upper bound is also specific to this construction
and these estimates; it rules out neither improved estimates nor a
different construction.

\section{The determinant construction}
\label{s2-construction}

This section follows Section~2 of the source, pp.~3--5.  We change the
functional to one adapted to $\zeta(9)$, while retaining the source's
matrix size, numerator $D_N^6$, and scalar normalization.  The degree and
positive integral representation can then be calculated directly.  The
arithmetically normalized decay statement is a separate conditional target.

\subsection{The functional and the matrix}
\label{s2-functional}

As in source (2.1), put
\begin{equation}
 K=40n,\qquad N=3n,\qquad h=K-N=37n,\qquad
 \alpha=\frac3{40},\quad\lambda=\frac{37}{40},\quad
 H=1+2\alpha=\frac{23}{20}.
 \label{s2-parameters}
\end{equation}
Define $D_0(t)=1$ and
\[
 D_m(t)=\prod_{j=1}^{m}(t+j^2),\qquad
 H_j^{(9)}=\sum_{v=1}^{j}\frac1{v^9},\qquad H_0^{(9)}=0.
\]
Use Bernoulli numbers with $B_1=-1/2$.  The analogues of source
(2.2)--(2.3) are
\begin{align}
 \mu_9(t^e)
 &=\frac{(-1)^e B_{2e+2}}{40320}
       \prod_{r=3}^{9}(2e+r),\qquad e\geq0,
 \label{s2-monomials}\\
 \mu_{9,X}\!\left(\frac1{t+j^2}\right)
 &=j^8\bigl(X-H_j^{(9)}\bigr)-\frac18+\frac1{2j},
       \qquad j\geq1.
 \label{s2-poles}
\end{align}
Here $40320=8!$, and $\mu_{9,X}$ agrees with $\mu_9$ on polynomials.
For example, $\mu_9(1)=3/4$ and $\mu_9(t)=11/8$; these provide simple
checks on the factorial.  Polynomial division and simple partial
fractions uniquely extend the functional to rational functions whose
poles are simple and lie among the negative integer squares.  Only the
pole contribution depends on $X$, and it does so affinely.
For instance, the first two pole values are
\[
 \mu_{9,X}\!\left(\frac1{t+1}\right)=X-\frac58,
 \qquad
 \mu_{9,X}\!\left(\frac1{t+4}\right)=256X-\frac{2051}{8}.
\]
The factor $256=2^8$ and the harmonic increment $2^{-9}$ must be changed
together.  Evaluating the old fifth-order functional at $X=\zeta(9)$
would retain its old moments and would not yield the weight constructed
below.  Both parts of the functional, not just its argument, are being
replaced.

Retaining source (2.4)--(2.5), define
\begin{equation}
 G_{9,K}(X)=
 \left[\mu_{9,X}\!\left(\frac{D_N(t)^6t^{i+j}}{D_K(t)}\right)
 \right]_{0\leq i,j<h},\qquad
 \Delta_{9,K}(X)=\det G_{9,K}(X),
 \label{s2-matrix}
\end{equation}
and
\begin{equation}
 S_K=\frac{(K!)^{2h}4^{h-1}}
 {(N!)^{12h}\displaystyle\prod_{i=1}^{h-1}((2i)!)^2},
 \qquad F_{9,K}(X)=S_K\Delta_{9,K}(X).
 \label{s2-normalization}
\end{equation}
The numerator remains $D_N^6$: the order of the target zeta value alone
does not require a different numerator power.  This choice also preserves
the leading real external field used later.  For fixed $M\geq40$, the
analogue of source (2.6) is
\begin{equation}
 Q_{9,K,M}(X)=m_{9,K,M}F_{9,K}(X),
 \qquad m_{9,K,M}\in\mathbb Q_{>0},
 \label{s2-Q}
\end{equation}
where Section~5 must supply the prime exponents and their cost.
Cancellation of $D_N$ leaves a denominator with distinct linear factors,
so all matrix entries lie in the stated domain.  Their polynomial
quotients can have substantial degree; they are retained in the
definition and cannot be discarded during the later denominator
calculation.  The factor $S_K$ is positive and independent of $X$.
Its reuse fixes the normalization being compared across the two zeta
orders, but by itself makes no assertion that every coefficient is
integral.

\subsection{The polynomial estimates: conditional targets}
\label{s2-targets}

\noindent\textbf{Draft analogue of Theorem 2.1 [Conditional].}
For fixed $M\geq40$ and $K\in40\mathbb Z_{>0}$ with $K\geq200M^2$,
the tasks for the later sections are to establish a positive rational
normalizer in \eqref{s2-Q} for which
\begin{equation}
 Q_{9,K,M}\in\mathbb Z[X],\qquad
 \deg Q_{9,K,M}=h,\qquad Q_{9,K,M}(\zeta(9))>0,
 \label{s2-integer-target}
\end{equation}
and to control its value by comparing the estimates
\begin{align}
 \log m_{9,K,M}&\leq A_{9,M}K^2+o_M(K^2),
 \label{s2-arithmetic-target}\\
 \log F_{9,K}(\zeta(9))&\leq U K^2+O(K\log K),
 \qquad U=-\frac{2733991}{2000000}.
 \label{s2-real-target}
\end{align}
If these estimates hold, their sum gives
\begin{equation}
 \log Q_{9,K,M}(\zeta(9))
 \leq (A_{9,M}+U)K^2+o_M(K^2).
 \label{s2-combined-target}
\end{equation}
A quadratic decay conclusion additionally requires $A_{9,M}+U<0$.
This sign is a goal, not an assumption or a conclusion of this section.
The numerical bounds in source (2.7)--(2.8) therefore have no asserted
$\zeta(9)$ counterparts here.  Degree and positivity will follow below
for every positive normalizer; coefficient integrality and the sign of
the combined estimate require the later arithmetic calculation.
In particular, a larger denominator-clearing multiplier can preserve
all three properties in \eqref{s2-integer-target} while destroying the
small-value estimate.  This is why the draft records its quadratic cost
as a separate quantity.  A favorable real integral estimate and exact
nonvanishing do not determine that cost.  Throughout this comparison,
$M$ is fixed before $K$ tends to infinity. The prime-sum argument in
Section~5 supplies a limiting quadratic bound, not an explicit
$O_M(K\log K)$ remainder for the arithmetic part.

\subsection{Degree}
\label{s2-degree}

\noindent\textbf{[Derived].}
Let $D_{\mathrm{tail}}=D_K/D_N$ and $W=D_N^5$.  Cancellation in a
matrix entry gives $Wt^{i+j}/D_{\mathrm{tail}}$.  At the pole $-k^2$,
its residue is
\[
 \frac{W(-k^2)(-k^2)^{i+j}}{D_{\mathrm{tail}}'(-k^2)}.
\]
Consequently, for the square Vandermonde matrix
$V_{ij}=(-j^2)^i$, indexed by $0\leq i<h$ and $N<j\leq K$,
\begin{equation}
 [X]G_{9,K}(X)=V\operatorname{diag}\!\left(
 \frac{j^8D_N(-j^2)^5}{D_{\mathrm{tail}}'(-j^2)}
 \right)_{N<j\leq K}V^{\mathsf T}.
 \label{s2-leading-matrix}
\end{equation}
There are exactly $h$ surviving pole locations.  If $x_j=-j^2$, then
\[
 \prod_jD_{\mathrm{tail}}'(x_j)
   =(-1)^{h(h-1)/2}\prod_{a<b}(x_b-x_a)^2.
\]
The Vandermonde square thus cancels the denominator in the determinant
of \eqref{s2-leading-matrix}.  The analogue of source (2.9) is
\begin{equation}
 [X^h]\Delta_{9,K}(X)
   =(-1)^{h(h-1)/2}
     \prod_{j=N+1}^{K}j^8D_N(-j^2)^5\ne0.
 \label{s2-leading-coefficient}
\end{equation}
Each entry is affine, so the determinant has degree at most $h$; this
nonzero coefficient makes its degree exactly $h$.  Multiplication by
$S_K$ or a positive rational $m_{9,K,M}$ preserves that degree.

\subsection{The positive integral representation}
\label{s2-positivity}

\noindent\textbf{Analogue of Proposition 2.2 [Derived].}
For every rational function in the domain of $\mu_{9,X}$,
\begin{equation}
 \mu_{9,\zeta(9)}(R)=\int_0^\infty R(y^2)w_9(y)\,dy,
 \qquad
 w_9(y)=\frac{(2\pi)^8y^9}{20160}
           \sum_{\ell\geq1}\ell^8e^{-2\pi\ell y}>0.
 \label{s2-integral}
\end{equation}
This specializes the general odd-order weight
$2(2\pi)^{s-1}y^s\sum_{\ell\geq1}\ell^{s-1}e^{-2\pi\ell y}/(s-1)!$
at $s=9$.  The specialization is useful because it changes the order
without introducing lower odd zeta values.

For monomials, integrating the nonnegative series term by term gives
\begin{align}
 \int_0^\infty y^{2e}w_9(y)\,dy
 &=\frac{(2e+9)!}{20160(2\pi)^{2e+2}}\zeta(2e+2)
 \notag\\
 &=\frac{(-1)^eB_{2e+2}}{40320}
       \frac{(2e+9)!}{(2e+2)!}
 =\mu_9(t^e).
 \label{s2-moment-integral}
\end{align}
The second equality uses Euler's even-zeta formula, as in source
Proposition~2.2.

For the pole calculation set $f(y)=(e^{2\pi y}-1)^{-1}$.  Since
$w_9(y)=y^9f^{(8)}(y)/20160$, the required derivative is obtained by
first dividing:
\[
 \frac{y^9}{y^2+a^2}
 =y^7-a^2y^5+a^4y^3-a^6y
    +a^8\frac{y}{y^2+a^2}.
\]
The polynomial part disappears after eight derivatives, while
$y/(y^2+a^2)=\operatorname{Im}(a-iy)^{-1}$.  Therefore
\begin{equation}
 \left(\frac{d}{dy}\right)^8\frac{y^9}{y^2+a^2}
   =8!\,a^8\operatorname{Im}(a-iy)^{-9}.
 \label{s2-eighth-derivative}
\end{equation}
Eight integrations by parts now give
\begin{equation}
 \int_0^\infty\frac{w_9(y)}{y^2+a^2}\,dy
   =2a^8\int_0^\infty
       \frac{\operatorname{Im}(a-iy)^{-9}}{e^{2\pi y}-1}\,dy.
 \label{s2-pole-ibp}
\end{equation}
To check the boundary terms, write $g(y)=y^9/(y^2+a^2)$.
For $0\leq r\leq7$, the product $g^{(7-r)}f^{(r)}$ is
$O(y^{r+2})O(y^{-r-1})=O(y)$ at zero.  At infinity the exponential
decay of every derivative of $f$ absorbs the polynomial factors.

Hermite's formula, the identity cited by the source as DLMF
(25.11.29), reads in this case
\[
 \zeta(9,a)=\frac1{2a^9}+\frac1{8a^8}
   +2\int_0^\infty
       \frac{\operatorname{Im}(a-iy)^{-9}}{e^{2\pi y}-1}\,dy,
 \qquad a>0.
\]
Thus \eqref{s2-pole-ibp} becomes
\begin{equation}
 \int_0^\infty\frac{w_9(y)}{y^2+a^2}\,dy
   =a^8\zeta(9,a)-\frac1{2a}-\frac18.
 \label{s2-hurwitz-pole}
\end{equation}
At $a=j\in\mathbb Z_{>0}$, substitute
$\zeta(9,j)=\zeta(9)-H_j^{(9)}+j^{-9}$ to recover exactly
\eqref{s2-poles}.  Together with the monomial calculation, linearity
proves \eqref{s2-integral}.  Moreover $w_9(y)\to1/\pi$ at zero and
$w_9(y)=O((1+y)^9e^{-2\pi y})$ at infinity, so every integral used
here converges.
The endpoint statements can also be read directly from the weight.
Near zero the Laurent expansion of $f$ gives
$f^{(8)}(y)=8!/(2\pi y^9)+O(1)$, which cancels the factor $y^9$.
For $y\geq1$,
\[
 \sum_{\ell\geq1}\ell^8e^{-2\pi\ell y}
 \leq e^{-2\pi y}
      \sum_{\ell\geq1}\ell^8e^{-2\pi(\ell-1)},
\]
and the final sum is a finite constant.  Thus rational functions in the
domain are bounded near the integration endpoint zero and have at most
polynomial growth at infinity, which the exponential decay absorbs.

Finally, if $q$ is a nonzero real polynomial of degree below $h$, then
\begin{equation}
 \int_0^\infty
   \frac{D_N(y^2)^6}{D_K(y^2)}q(y^2)^2w_9(y)\,dy>0.
 \label{s2-positive-quadratic-form}
\end{equation}
This is the quadratic form of $G_{9,K}(\zeta(9))$ at the coefficient
vector of $q$.  All weight factors are positive for $y>0$, and $q(y^2)$
has only finitely many zeros.  Hence the matrix is positive definite,
$\Delta_{9,K}(\zeta(9))>0$, and every positive normalization retains
nonvanishing.  This completes the local algebraic construction without
assuming the later coefficient-clearing cost is small enough.

\section{The local rational functional}
\label{s3-functional}

This section translates the rational-functional calculation in Fauzan,
Section~3, pp.~5--9, including (3.1)--(3.12).  The parameter $X$ remains
indeterminate: a valuation estimate must control every coefficient, including
the constant coefficient.  Write $v_p^G(A)$ for the minimum valuation of the
coefficients of $A\in\mathbb Q_p[X]$, with $v_p^G(0)=+\infty$.

\paragraph{\textnormal{[Derived] Pullback, reflection, and difference.}}
Let $L(x^d)=B_d$, with $B_1=-1/2$, and define
\begin{equation}
 \tau_9(P)=\frac{L(P^{(7)})}{8!},\qquad
 \tau_{9,X}\!\left(\frac1{x-r}\right)=H_{d(r)}^{(9)}-X,
 \qquad
 d(r)=\begin{cases}r,&r\geq0,\\-r-1,&r<0.\end{cases}
 \label{s3-tau}
\end{equation}
Extend by polynomial division and simple partial fractions.  The analogue
of (3.1) is
\begin{equation}
 \mu_{9,X}(R)=\tau_{9,X}\bigl(x^9R(-x^2)\bigr).
 \label{s3-pullback}
\end{equation}
Indeed, seven differentiations of $(-1)^e x^{2e+9}$ give
$(-1)^e B_{2e+2}\prod_{r=3}^{9}(2e+r)/8!$.  At a pole, division gives
\[
 \frac{x^9}{j^2-x^2}
 =-x^7-j^2x^5-j^4x^3-j^6x
  -\frac{j^8}{2}\left(\frac1{x-j}+\frac1{x+j}\right).
\]
Only $-x^7$ contributes to the polynomial functional, with value $-1/8$.
Using $H_j^{(9)}-H_{j-1}^{(9)}=j^{-9}$ therefore recovers
$j^8(X-H_j^{(9)})-1/8+1/(2j)$, including the easily missed constant term.

The Bernoulli identities $L(P(-1-x))=L(P(x))$ and
$L(P(x+1)-P(x))=P'(0)$ yield the analogues of (3.2)--(3.3):
\begin{align}
 \tau_{9,X}(g(-1-x))&=-\tau_{9,X}(g(x)),
 \label{s3-reflection}\\
 \tau_{9,X}(g(x+1)-g(x))&=\frac{g^{(8)}(0)}{8!}.
 \label{s3-difference}
\end{align}
The second identity requires $g$ to be regular at zero.  For integer poles,
reflection uses $d(-r-1)=d(r)$; the difference identity uses the same
harmonic increment, whose value for $g=1/(x-r)$ is $-r^{-9}$.

\subsection{Completion at a prime}

\paragraph{\textnormal{[Derived] Analytic and principal-part pieces.}}
Fix $p\geq11$.  Put
\[
 \kappa_d=\frac{d(d-1)\cdots(d-6)}{8!}B_{d-7}\quad(d\geq7),
 \qquad \kappa_0=\cdots=\kappa_6=0.
\]
Since $p\nmid8!$, von Staudt--Clausen gives $v_p(\kappa_d)\geq-1$.
Also $\kappa_d\in\mathbb Z_p$ for $d\leq p+5$, because then the relevant
Bernoulli index is at most $p-2$.  This is a sufficient degree cutoff;
no optimal cutoff is needed here.  Consequently
\[
 \tau_9\!\left(\sum_{d\geq0}f_dz^d\right)=\sum_{d\geq0}f_d\kappa_d,
 \qquad \|\tau_9\|\leq p,
 \qquad \mathcal A=\mathbb Q_p\langle z\rangle .
\]
Thus the analytic part can lose one power of $p$.

Let $T(z)=\prod_\nu(z-r_\nu)$, where the integer roots have pairwise
unit differences and $d(r_\nu)<p$.  Polynomial division and partial
fractions have norm at most one, since
$c_\nu=U(r_\nu)/T'(r_\nu)$ and $T'(r_\nu)$ is a unit.  They extend to
the direct sum
\[
 \mathcal B_T=\mathcal A\oplus\bigoplus_\nu
       \mathbb Q_p\frac1{z-r_\nu},\qquad
 \tau^{\mathrm{ext}}_{9,Y}\!\left(f+\sum_\nu\frac{c_\nu}{z-r_\nu}\right)
 =\tau_9(f)+\sum_\nu c_\nu(H_{d(r_\nu)}^{(9)}-Y).
\]
For $Y\in\mathbb Z_p[X]$, the principal-part values are integral.
Accordingly, the analogue of (3.4) is
\begin{equation}
 \tau^{\mathrm{ext}}_{9,Y}
 \!\left(\sum_{j\geq0}\frac{p^jU_j(z)}{T(z)}\right)
 \in\mathbb Z_p[X],\qquad
 U_j\in\mathbb Z_p[z],\quad\deg U_0\leq p+5.
 \label{s3-completion}
\end{equation}
The initial quotient has integral functional value by its degree; every
later term carries the extra $p$ needed to absorb the analytic loss.
Convergence follows in the maximum norm on $\mathcal B_T$.

For a far pole, $v_p(s)<0$, expand $1/(z-s)=-\sum_{d\geq0}z^ds^{-d-1}$.
Applying $\tau_9$ termwise gives the explicit analogue of (3.5):
\begin{equation}
 \tau_9^{\mathrm{an}}\!\left(\frac1{z-s}\right)
 =-\frac18\sum_{k\geq0}\binom{k+7}{7}B_k\,s^{-k-8}.
 \label{s3-far-pole}
\end{equation}
The series converges because $v_p(s^{-1})>0$ and the Bernoulli valuations
are bounded below.  The decomposition into analytic functions and the
specified principal parts is unique.  Hence expanding far factors first
and taking global partial fractions first produce the same value.

\subsection{Distribution}

\paragraph{\textnormal{[Derived] The new scale is $p^{-8}$.}}
Define the analogues of the constants in (3.6) by
\begin{equation}
 C_{p,9}=\sum_{a=1}^{p-1}\tau_9^{\mathrm{an}}
                  \!\left(\frac1{z+a/p}\right),
 \qquad Y=p^9X+C_{p,9}.
 \label{s3-distribution-constant}
\end{equation}
Here $C_{p,9}\in p^9\mathbb Z_p$.  In fact, the $k=0$ part of
\eqref{s3-far-pole} is $-p^8\sum_{a=1}^{p-1}a^{-8}/8$, and its residue
vanishes because $p-1\nmid8$ for $p\geq11$.  The $k=1$ part already
has valuation at least $9$; for $k\geq2$ the bound is $k+7\geq9$.
The condition $p\geq11$ also makes $8!$ a unit.  Smaller primes are
handled separately below.

The analogue of (3.7), for rational $g$ with simple integer poles, is
\begin{equation}
 \tau_{9,X}(g)=p^{-8}\sum_{a=0}^{p-1}
        \tau^{\mathrm{ext}}_{9,Y}\bigl(g(a+pz)\bigr).
 \label{s3-distribution}
\end{equation}
Bernoulli multiplication contributes $p^{-1}$ and the seven derivatives
contribute $p^{-7}$, proving the polynomial identity.  At $g=1/x$ the
right side is $p^{-9}(-Y+C_{p,9})=-X$, which explains the parameter choice.
The substitutions $a\mapsto p-1-a$, $z\mapsto-1-z$ preserve reflection.
Translation telescopes to
\[
 p^{-8}\tau^{\mathrm{ext}}_{9,Y}
       \bigl(g(p(z+1))-g(pz)\bigr)=\frac{g^{(8)}(0)}{8!}.
\]
The analytic identities extend by continuity.  Starting with the pole
at zero, translation determines positive integer poles and reflection
determines negative ones.  This establishes the whole rational identity,
including its polynomial contribution.

\subsection{Small primes and the integer-valued basis}

\paragraph{\textnormal{[Derived] A uniform, deliberately coarse bound.}}
For any prime $p$, an integer-valued polynomial $A$ of degree at most $d$
has integral coefficients in the basis $\binom{x}{k}$, given by forward
differences.  Vandermonde's identity gives the unchanged analogue of (3.8):
\[
 v_p(A(x)-A(y))\geq v_p(x-y)-\lfloor\log_p\max(1,d)\rfloor.
\]
Its $r$th derivative loses at most $r\lfloor\log_p\max(1,d)\rfloor$:
the coefficient of $z^r$ in $\binom{z}{j}$ is a sum of products of $r$
reciprocals of integers at most $j$.  Moreover
$L(\binom{x}{k})=(-1)^k/(k+1)$, from
$L((1+z)^x)=\log(1+z)/z$.  Seven derivatives followed by $L$ therefore
give the analogue of (3.9):
\begin{equation}
 P(\mathbb Z_p)\subset p^{-A}\mathbb Z_p,\quad\deg P\leq d
 \quad\Longrightarrow\quad
 v_p(\tau_9(P))\geq-A-8\lfloor\log_p(d+1)\rfloor-v_p(8!).
 \label{s3-small-polynomial}
\end{equation}

For integer-valued $A$, $\deg A\leq d$, set
\[
 g(x)=\frac{(K!)^2A(x)}{\prod_{-K\leq r\leq K,\ r\ne0}(x-r)},
 \qquad L_0=\lfloor\log_p(2K)\rfloor,
 \qquad M=\lfloor\log_p\max(2K,d+1)\rfloor.
\]
Its residue at $r$ is
$c_r=(-1)^{K-r}(K!)^2rA(r)/((K+r)!(K-r)!)$, of valuation at least
$-L_0$.  The root-counting and inverse-product estimates of source
Lemma~3.3 are independent of the zeta index.  Outside the balls
$r+p^{L_0+1}\mathbb Z_p$ they bound the rational terms by $-2L_0$;
inside a ball, subtracting its principal part and using the preceding
difference estimate bounds the polynomial quotient $P$ by $-2M$.
Applying \eqref{s3-small-polynomial} gives a loss at most $10M+v_p(8!)$.
The pole terms have the same bound, since
$v_p(H_{d(r)}^{(9)})\geq-9\lfloor\log_pK\rfloor$.
Thus the analogue of (3.10) is
\begin{equation}
 v_p^G(\tau_{9,X}(g))\geq-10M-v_p(8!).
 \label{s3-small-rational}
\end{equation}

Keep exactly the source basis and normalization:
\[
 q_0=1,\qquad q_i(t)=\frac{(-1)^i2tD_{i-1}(t)}{(2i)!}\ (i\geq1),
 \qquad f_i(t)=\left(\frac{D_N(t)}{(N!)^2}\right)^3q_i(t).
\]
Both identities
\[
 q_i(-x^2)=\binom{x+i}{2i}+\binom{x+i-1}{2i},\qquad
 \frac{D_N(-x^2)}{(N!)^2}=\binom{N-x}{N}\binom{N+x}{N}
\]
show that $f_i(-x^2)$ is integer-valued on $\mathbb Z_p$.
Squaring the determinant of this triangular basis change gives precisely
the prescribed $S_K$.  Hence the analogue of (3.11) is
\begin{equation}
 F_{9,K}(X)=\det\left[
 (K!)^2\mu_{9,X}\!\left(\frac{f_i(t)f_j(t)}{D_K(t)}\right)
 \right]_{0\leq i,j<h}.
 \label{s3-normalized-basis}
\end{equation}
The pullback numerator, up to the harmless sign $(-1)^K$, is
$x^9f_i(-x^2)f_j(-x^2)$.  Its degree is at most
$12N+4h+5=23K/5+5$.  For $K\in40\mathbb Z_{>0}$ this gives
$\max(2K,d+1)\leq5K$.  Entrywise use of
\eqref{s3-small-rational} and the determinant expansion finally yield
the analogue of (3.12):
\begin{equation}
 \boxed{\displaystyle
 v_p^G(F_{9,K})\geq
 -10h\lfloor\log_p(5K)\rfloor-hv_p(8!)
 \qquad\text{for every prime }p.}
 \label{s3-small-final}
\end{equation}
This accounts for small-prime denominators.  Whether the sharper
large-prime estimates leave enough real decay is a separate question.

\section{Local estimates for the determinant}
\label{s4-local}

This section follows Section~4 of the source, pp.~9--13, with
$D_N^6$, $h=K-N$, and the scalar $S_K$ unchanged.  Write $v_p^G$ for
the minimum valuation of the coefficients of a polynomial in $X$.
The object estimated here is $\Delta_{9,K}$ before multiplication by
$S_K$.  Fix $M\geq40$.  The inner and outer bases serve different
purposes: the first spreads vanishing among residue classes, while the
second cancels the poles whose harmonic sums contain multiples of $p$.

\subsection{The inner range and its new row weights}
\label{s4-inner}

Assume, as in source (4.1), that
\begin{equation}
 K\in40\mathbb Z_{>0},\qquad K\geq200M^2,
 \qquad K/M<p\leq K/3.
 \label{s4-inner-range}
\end{equation}
Then $p>200M$ and $p^2>5K$.  Put
\[
 m_A=\lfloor A/p\rfloor,\qquad
 \ell_A(a)=\#\{1\leq j\leq A:j\equiv\pm a\pmod p\},
 \qquad 1\leq a\leq m=(p-1)/2.
\]
If a row polynomial vanishes to order $\nu_i(a)$ at $t=-a^2$,
the ordinary-class bound replacing source (4.2) is
\begin{equation}
 \nu_i(a)+\nu_j(a)+6\ell_N(a)-\ell_K(a)-8.
 \label{s4-ordinary-source}
\end{equation}
At zero the bound replacing (4.3) is unchanged:
\begin{equation}
 2\nu_i(0)+2\nu_j(0)+12m_N-2m_K+1.
 \label{s4-zero-source}
\end{equation}
\textbf{[Derived]}
Indeed, the distribution factor is now $p^{-8}$.  At an ordinary
source $x=a+pz$, $x^9$ is a unit to leading order, giving the loss
$-8$.  At $x=pz$, its factor $p^9$ compensates for that distribution
factor and leaves $+1$.
\textbf{[Formal analogue]} Applying these counts to the entire rational
summand, including its polynomial quotient, requires the analytic
completion estimate from Section~3.  We keep that transplant explicit.

Retain the source allocation, namely
\begin{equation}
 \begin{gathered}
 L_0=4M+10,\qquad b_a=3\ell_N(a),\\
 mT+E=h-L_0+3(N-m_N),\qquad 0\leq E<m,\\
 L_a=T-b_a+\epsilon_a,\qquad Z_a=L_a+b_a=T+\epsilon_a.
 \end{gathered}
 \label{s4-allocation}
\end{equation}
Here $\epsilon_a=1$ for the first $E$ classes in decreasing order of
$\ell_K(a)$, and is zero otherwise.  With $x=K/p$, the unchanged
inequalities
$2Hx-21/20<T<2Hx$ and $b_a\leq6\alpha x+3$ give
$T-b_a>2\lambda x-81/20\geq3/2$.  Thus the dimensions are
nonnegative and sum to $h$.

Use the polynomials from source (4.5):
\begin{equation}
 E_{a,i}(t)=\prod_{\substack{0\leq c\leq m\\c\ne a}}
 (t+c^2)^{L_c}(t+a^2)^i,
 \qquad 0\leq i<L_a.
 \label{s4-inner-basis}
\end{equation}
Their degrees are less than $h$.  Distinct factors are coprime modulo
$p$, and their local Taylor-coordinate matrices are triangular with
unit diagonals.  Hence this is a $\mathbb Z_p$-unimodular basis.
The weights replacing source (4.6)--(4.8) are
\begin{align}
 w^{(9)}_{a,i}&=i+b_a-\frac{\ell_K(a)+8}{2},\qquad a>0,
 \label{s4-inner-weights}\\
 w^{(9)}_{0,i}&=\min\left\{
 2i+6m_N-m_K+\frac12,
 \min_{1\leq c\leq m}\left(Z_c-\frac{\ell_K(c)+8}{2}\right)
 \right\},
 \label{s4-inner-zero-weight}\\
 \gamma^{\mathrm{in}}_{9,p}&=
 2\sum_{a=0}^{m}\sum_{i=0}^{L_a-1}w^{(9)}_{a,i}.
 \label{s4-inner-gamma}
\end{align}

The comparison between two ordinary blocks reduces, just as before,
to
$Z_c-Z_a+1+(\ell_K(a)-\ell_K(c))/2\geq0$.
The counts take two consecutive values, and the ordering of the
extras resolves the only potentially adverse case.  At zero the
half-weight of an ordinary row is at least $7M+41/2$, comfortably
above its assigned weight.  Formula~\eqref{s4-inner-zero-weight}
handles the zero-block rows at every source.

The local degree checks also survive.  At an ordinary source the
leading numerator degree is at most $2(T+1)<23M/5+2$; at zero it is
at most
\[
 9+4L_0+12m_N\leq169M/10+49<p+1.
\]
\textbf{[Conditional]}
Subject to the Section~3 analytic completion estimate, every matrix
entry in this basis has valuation at least the sum of its row
weights.  Expanding the determinant therefore gives
\begin{equation}
 v_p^G(\Delta_{9,K})\geq\gamma^{\mathrm{in}}_{9,p}.
 \label{s4-inner-bound}
\end{equation}
Relative to the source weights, every ordinary weight decreases by
two; every zero-block weight decreases by an amount between zero
and two.  Consequently the useful exact comparison is
\begin{equation}
 -4h\leq\gamma^{\mathrm{in}}_{9,p}-\gamma^{\mathrm{in}}_{5,p}
 \leq-4h+4L_0.
 \label{s4-inner-shift}
\end{equation}
Thus the limiting inner cost numerator changes by $4\lambda x$,
with only an $O_M(1)$ error in the finite exponent.

\subsection{Removing the outer polynomial correction}
\label{s4-rank-correction}

Use the source outer hypotheses, with the harmless strengthened
threshold $p\geq11$:
\begin{equation}
 p\leq K<3p,\qquad p^2>2K,\qquad 2N<p,
 \qquad5N\leq2p-2.
 \label{s4-outer-range}
\end{equation}
\textbf{[Derived]}
The denominator $40320$ is a $p$-adic unit.  Von Staudt--Clausen
therefore gives $v_p(\mu_9(t^e))\geq-1$.
If $e<4p-5$ and a Bernoulli denominator contains $p$, then
$2e+2=j(p-1)$ for some $1\leq j\leq7$.
The factor $2e+2+j=jp$ occurs among $2e+3,\ldots,2e+9$ and
cancels that denominator.  Hence
\begin{equation}
 \mu_9(t^e)\in\mathbb Z_p\qquad(e<4p-5).
 \label{s4-moment-cutoff}
\end{equation}

Choose $c_e\in\mathbb Z_p$ representing $p\mu_9(t^e)$ modulo $p$,
with $c_e=0$ below this cutoff.  Replace the polynomial moments by
$\mu_9^0(t^e)=\mu_9(t^e)-c_e/p$, leaving simple-pole values
unchanged.  Set $W=D_N^5$ and $D_{\mathrm{tail}}=D_K/D_N$;
the rational input is still $Wt^{i+j}/D_{\mathrm{tail}}$.
Its polynomial quotient has degree at most $i+j+6N-K$.  Therefore
\begin{equation}
 G_{9,K}=A+p^{-1}L,\qquad
 \operatorname{rank}_{\mathbb Q_p}L\leq
 r_{9,p}:=\max(0,K+4N-4p+4),
 \label{s4-matrix-rank}
\end{equation}
where $L$ is integral.  Indeed,
$L_{ij}=0$ if $i+j<K-6N+4p-5$, so its first $h-r_{9,p}$
rows and columns vanish exactly.  This is a rank bound over
$\mathbb Q_p$, rather than only after reduction modulo $p$.
For our fixed parameters and $p>K/3$,
$K+4N-4p+4<4-K/30$; thus $r_{9,p}=0$ throughout the outer
range once $K\geq120$.

\subsection{The outer basis, class costs, and rank lemma}
\label{s4-outer-basis}

For an ordinary square class put $\ell=\ell_K(a)$ and
$\delta=\mathbf1_{a\leq N}$.  Let $Q_a$ contain its remaining
$\ell-\delta$ factors, $P_a=D_{\mathrm{tail}}/Q_a$, and let
$E_a$ contain its $\ell-2$ factors with $j>p$.  The rows are
$P_aq_{a,i}$, with
\begin{equation}
 q_{a,i}(t)=
 \begin{cases}
 (t+a^2)^i,&0\leq i<\ell-2,\\
 E_a(t)(t+a^2)^{i-(\ell-2)},&\ell-2\leq i<\ell-\delta.
 \end{cases}
 \label{s4-outer-polynomials}
\end{equation}
Together with the zero-class rows these form a unimodular basis:
each local sequence is monic of successive degrees, and distinct
classes have unit resultants.  Cross-class entries of $A$ are
integral because their quotients are integral polynomials.

\textbf{[Derived]}
For two low indices, the row factors supply $i+j$, $W$ supplies
$5\delta$, the node derivative loses $\ell-\delta-1$, and the
harmonic value loses at most nine powers of $p$.  The resulting
bound is $i+j+6\delta-\ell-8$.  A high-index row cancels all
poles with $j>p$.  The remaining values at $a,p-a$ are integral
and congruent modulo $p$: use
$H_{p-a}^{(9)}\equiv H_{a-1}^{(9)}\pmod p$ and the displayed
simple-pole formula for $\mu_{9,X}$.  The $X$ coefficients agree
because $(p-a)^8\equiv a^8$.  The constant-coefficient difference
also vanishes: the extra $a^{-9}$ in $H_a^{(9)}$ cancels the
change in $1/(2a)$.  The divided difference is consequently
integral, including after multiplication by polynomial numerators.

Thus the analogue of source (4.12) is
\begin{equation}
 w^{(9)}_{a,i}=
 \begin{cases}
 \min\{0,i+3\delta-(\ell+8)/2\},&i<\ell-2,\\
 0,&i\geq\ell-2.
 \end{cases}
 \label{s4-outer-weights}
\end{equation}
At zero a single pole has weight $-1/2$; for $p,2p$, rows
$1,t+p^2$ give entry bounds $-3,-1,0$, so weights $-2,0$
remain valid.  Here each pole value has valuation at least $-1$
because $j^8H_j^{(9)}$ loses only one power when $p\mid j$.
The spacing of the two squared nodes costs two further powers.

For reference, the elementary rank lemma from source (4.13)
is unchanged.  If $v(A_{ij})\geq w_i+w_j$, the $w_i$ are
nonpositive half-integers, $L$ is integral of rank at most $r$,
and exactly $z$ weights are zero, then
\begin{equation}
 v\det(A+p^{-1}L)\geq2\sum_iw_i-\min(r,z).
 \label{s4-rank-lemma}
\end{equation}
A complementary-minor term using $k$ rows and columns of $L$
has valuation at least
$2\sum_iw_i-k-\sum_{i\in I}w_i-\sum_{j\in J}w_j$.
Terms with $k>r$ vanish.  Ordering weights decreasingly, the
largest possible loss is
$\max_{0\leq k\leq\min(r,h)}(k+2\sum_{i=1}^kw_{(i)})$:
each zero weight increases it by one, and every subsequent
increment is nonpositive.  This proves the lemma, also for
Gauss valuations.  Unimodular changes of basis preserve the
integrality and rank of $L$.

\subsection{Explicit outer exponent and the primes beyond \texorpdfstring{$K$}{K}}
\label{s4-outer-formula}

\textbf{[Derived]}
For a removed class, the new table is
\[
\begin{array}{c|rrrrr}
 \ell&2&3&4&5&6\\\hline
 -2\sum_iw^{(9)}_{a,i}&0&5&10&15&20\\
 \#\{i:w^{(9)}_{a,i}=0\}&1&1&1&1&1
\end{array}
\]
An unremoved class costs $11(\ell-2)$ and has two zero weights.
Put
\[
 q=\lfloor K/p\rfloor\in\{1,2\},\quad v=K-qp,\quad
 u=\max(0,N+v-p+1),\quad t_p=\min(N,v)+u.
\]
Over all ordinary classes,
$\sum(\ell-2)=K-p-q+1$.  Over the $N$ removed classes it is
$2(q-1)N+t_p$: the two indicators $a\leq v$ and $a\geq p-v$
account for $t_p$.  Each removed class saves $6(\ell-2)$ from
the unremoved cost.  Adding the zero-class costs $1,4$ yields
\begin{equation}
 \gamma^{\mathrm{out}}_{9,p}=
 \begin{cases}
 -11(K-p)+6t_p-1-\min(r_{9,p},p-1-N),&K<2p,\\
 -11(K-p)+7+12N+6t_p-\min(r_{9,p},p-N),&K\geq2p.
 \end{cases}
 \label{s4-outer-gamma}
\end{equation}
The zero counts are $p-1-N$ and $p-N$, respectively.  Unlike
the source table, the removed classes never acquire additional
zero weights as $\ell$ increases.  Applying
\eqref{s4-rank-lemma} gives
$v_p^G(\Delta_{9,K})\geq\gamma^{\mathrm{out}}_{9,p}$.
As a finite consistency check, $K=40,N=3,p=29$ gives $-104$,
matching the independently computed constant-coefficient
valuation; this check does not establish the general estimate.

For $p>K$, every pole is below $p$ and each square class has at
most two poles.  The same congruence and divided-difference
argument makes its contribution integral.  The largest possible
polynomial quotient degree is $K+4N-2<4p-5$, so no polynomial
correction remains.  Consequently
\begin{equation}
 v_p^G(\Delta_{9,K})\geq0\qquad(p>K).
 \label{s4-large-prime-integrality}
\end{equation}
These local statements supply normalization exponents.  They do
not by themselves supply the final decay inequality.

\section{Normalization and the prime sum}
\label{s5-normalization}

We keep $K=40n$, $N=3n$, $h=37n$, $\alpha=3/40$, and
$\lambda=37/40$.  In particular, the scalar $S_K$ is unchanged.
The purpose of normalization is to turn the local determinant estimates
into an integral polynomial while recording the entire cost, including
positive valuations which permit division by a prime power.

\paragraph{[Conditional] The normalization exponents.}
The analogue of source (5.1)--(5.3), for fixed $M\geq40$ and
$K\geq200M^2$, is
\begin{equation}
 L_{9,p}(K,M)=
 \begin{cases}
 -10h\lfloor\log_p(5K)\rfloor-hv_p(40320),&pM\leq K,\\
 v_p(S_K)+\gamma^{\mathrm{in}}_{9,p},&pM>K,\ 3p\leq K,\\
 v_p(S_K)+\gamma^{\mathrm{out}}_{9,p},&3p>K.
 \end{cases}
 \label{s5-local-exponents}
\end{equation}
Here the two $\gamma$'s are the explicit weight sums and outer formulas
of Section~4, with $\gamma^{\mathrm{out}}_{9,p}=0$ for $p>K$.  Set
\begin{equation}
 m_{9,K,M}=\prod_{p\leq2h}p^{-L_{9,p}(K,M)},\qquad
 Q_{9,K,M}=m_{9,K,M}F_{9,K}.
 \label{s5-normalizing-product}
\end{equation}
No positive-part operation is applied to $-L_{9,p}$: the product is a
positive rational number, not necessarily an integer.  Legendre's formula gives
\begin{align}
 v_p(S_K)={}&2h\sum_{a\geq1}\lfloor K/p^a\rfloor
 -12h\sum_{a\geq1}\lfloor N/p^a\rfloor\notag\\
 &-2\sum_{i=1}^{h-1}\sum_{a\geq1}\lfloor2i/p^a\rfloor
 +(h-1)v_p(4).
 \label{s5-legendre}
\end{align}
Subject to the local estimates identified in Sections~3--4,
$Q_{9,K,M}\in\mathbb Z[X]$, the analogue of Proposition~5.1.
Indeed, these estimates handle $p\leq2h$; for $p>2h$ every factorial
in $S_K$ is a unit and the determinant is integral.  The analytic
completion step needed for the inner estimate remains a stated
dependency, rather than a theorem established by this normalization.

\subsection{The inner limiting function}
\label{s5-inner}

\paragraph{[Derived] The change from $5$ to $9$.}
For $x\geq3$ and $0<z<1/2$, put
\[
 \ell(x,z)=\lfloor x-z\rfloor+\lfloor x+z\rfloor+1,
 \qquad b(x,z)=3\ell(\alpha x,z).
\]
With $H=23/20$, define
$T=\lfloor2Hx\rfloor$, $s=Hx-T/2$, $q=\lfloor2x\rfloor$,
and $n_+=(2x-q)/2$.  The limiting weight sum is
\begin{align}
 \Gamma_9(x)={}&\int_0^{1/2}(T-b(x,z))
       (T+b(x,z)-\ell(x,z)-9)\,dz\notag\\
 &+s(2T-q-9)+(s-n_+)_+.
 \label{s5-gamma9}
\end{align}
Thus the constant $5$ in source (5.4) becomes $9$.  The allocation of
rows, including the last positive-part term, stays the same.  Since
$\int_0^{1/2}b(x,z)\,dz=3\alpha x$ and $H-3\alpha=\lambda$,
subtracting the source expression gives exactly
\[
 \Gamma_9(x)-\Gamma_5(x)
 =-4\left(T/2-3\alpha x+s\right)=-4\lambda x.
\]
The scalar still has limiting function
\begin{gather}
 J(u)=mu-\frac{m(m+1)}4,\quad m=\lfloor2u\rfloor,\qquad
 \mathcal N(x)=2\lambda x\lfloor x\rfloor
 -12\lambda x\lfloor\alpha x\rfloor-2J(\lambda x),\notag\\
 R_9(x)=-\Gamma_9(x)-\mathcal N(x)=R_5(x)+4\lambda x.
 \label{s5-r9-shift}
\end{gather}
Consequently, the analogue of (5.7) is
\[
 \gamma^{\mathrm{in}}_{9,p}=p\Gamma_9(K/p)+O_M(1),\qquad
 v_p(S_K)=p\mathcal N(K/p)+O_M(1).
\]
The finitely many reserved zero-block rows affect only the error term.
This is an algebraic comparison of the certificates; applying them to
the actual determinant still needs the local input.

\subsection{The outer limiting function}
\label{s5-outer}

\paragraph{[Derived] A vanishing rank correction.}
Write $y=p/K$.  In this range $y>1/3$, so
$(1+4\alpha-4y)_+=0$.  More concretely,
$r_{9,p}=\max(0,K+4N-4p+4)=0$ for all sufficiently large $K$.
The new raw determinant cost is therefore
\begin{equation}
 C_9(y)=
 \begin{cases}
 11(1-y)-12\alpha-6\min(\alpha,1-2y)
       -6(1+\alpha-3y)_+,&1/3<y<1/2,\\
 11(1-y)-6\min(\alpha,1-y)
       -6(1+\alpha-2y)_+,&1/2<y<1,\\
 0,&y>1.
 \end{cases}
 \label{s5-c9}
\end{equation}
For example, when $K\geq2p$, the limits of $v/K$ and $u/K$ are
$1-2y$ and $(1+\alpha-3y)_+$; inserting these in Section~4 gives
the first line.  This replaces both $R_0-d$ in source (5.8)--(5.9).
There is no additional limiting rank subtraction.

The scalar contribution is unchanged:
\[
 S(y)=-2\lambda\lfloor1/y\rfloor+
              \sum_{j=1}^{5}(2\lambda-jy)_+.
\]
Exact integration, detailed in Appendix~B, gives
\begin{gather}
 \int_{1/3}^{1}C_9(y)\,dy=\frac{14441}{7200},\qquad
 \int_{1/3}^{2\lambda}S(y)\,dy=\frac{1851}{16000},\notag\\
 I_{\mathrm{out},9}=\int_{1/3}^{2\lambda}(C_9(y)+S(y))\,dy
                   =\frac{305479}{144000}.
 \label{s5-iout9}
\end{gather}
Relative to the source's $I_{\mathrm{out},5}=127751/96000$, the
increase is $45541/57600$.  This calculation includes the improvement
in polynomial-moment integrality; it does not simply multiply every
old local loss by a common factor.

\paragraph{[Conditional] Summing over primes.}
The analogue of Proposition~5.2 is
\begin{equation}
 \limsup_{\substack{K\to\infty\\40\mid K}}
 \frac{\log m_{9,K,M}}{K^2}
 \leq I_{\mathrm{out},9}+\frac{10\lambda}{M}
                    +\int_3^M\frac{R_9(x)}{x^3}\,dx.
 \label{s5-prime-sum}
\end{equation}
Primes $p\leq\sqrt{5K}$ contribute $O(K^{3/2}\log K)$; this also
absorbs the fixed exceptional factors in $40320$.  For larger primes
in the small range, the logarithmic factor is one, giving
$10\lambda/M$ by the prime number theorem.  For bounded piecewise
continuous $f$, partial summation gives
\[
 K^{-2}\sum_{K/M<p\leq K/3}p f(K/p)\log p
       \longrightarrow\int_3^M f(x)x^{-3}\,dx.
\]
The $O_M(1)$ local errors contribute $o(K^2)$, and the same change of
variables in the outer range gives \eqref{s5-iout9}.

\subsection{The tail integral}
\label{s5-tail}

\paragraph{[Derived] The periodic part is unchanged.}
Put $f=\{x\}$, $g=\{\alpha x\}$, $\tau=\{2Hx\}$,
$\sigma=\{2x\}$, $\eta=\{2\lambda x\}$, and
$A(z)=\ell(x,z)-2x$, $B(z)=\ell(\alpha x,z)-2\alpha x$.
Expanding \eqref{s5-gamma9} gives the analogue of (5.12)--(5.14):
\begin{align}
 R_9(x)&=xF_9(x)+Q(x),\qquad
 F_9(x)=8\lambda+2\lambda f-12\lambda g,\notag\\
 Q(x)&=\frac{\tau(\tau-\sigma)-(\tau-\sigma)_+
                     +\eta(1-\eta)}2
       +9\int_0^{1/2}B^2\,dz-3\int_0^{1/2}AB\,dz,
 \label{s5-tail-decomposition}\\
 -\frac12&\leq Q(x)\leq\frac{13}{8}.\notag
\end{align}
Only the mean linear term has changed.  Define the same periodic primitives
\begin{gather*}
 P(x)=74g(1-g)-\lambda f(1-f),\qquad \overline P=\frac{2923}{240},\\
 G_0(v)=\frac{v(1-v)(2v-1)}6,\qquad
 C_{\mathrm{per}}(x)=\frac{74}{\alpha}G_0(g)-\lambda G_0(f).
\end{gather*}
Then $P'=F_9-3\lambda$, $C_{\mathrm{per}}'=P-\overline P$, and
$|C_{\mathrm{per}}|<16$.  Thus source (5.15) becomes
\begin{align}
 \int_T^\infty\frac{R_9(x)}{x^3}\,dx
 ={}&\frac{3\lambda}{T}-\frac{P(T)}{T^2}
       +\frac{\overline P}{T^2}
       -\frac{2C_{\mathrm{per}}(T)}{T^3}\notag\\
 &+6\int_T^\infty\frac{C_{\mathrm{per}}(x)}{x^4}\,dx
       +\int_T^\infty\frac{Q(x)}{x^3}\,dx.
 \label{s5-tail-identity}
\end{align}
At $T=20$, $P(T)=37/2$ and $C_{\mathrm{per}}(T)=0$.  At $40\mid M$,
both periodic primitives vanish.  Accordingly,
\begin{align}
 \int_{20}^\infty R_9(x)x^{-3}\,dx&\leq\frac{6191}{48000},
 \label{s5-tail-upper}\\
 \int_M^\infty R_9(x)x^{-3}\,dx&\geq\frac{3\lambda}{M}
       +\frac{2923/240-1/4}{M^2}-\frac{32}{M^3}.
 \label{s5-tail-lower}
\end{align}
Appendix~B evaluates the finite integral as
\begin{equation}
 \int_3^{20}R_9(x)x^{-3}\,dx
 =\frac{8222332206149897772829}{7535670527041937280000}.
 \label{s5-finite-inner}
\end{equation}
Combining these expressions yields the analogues of (5.19)--(5.21):
\begin{align}
 A_{9,*}&=\frac{25180312106351244019069}{7535670527041937280000},
 \notag\\
 A_{9,M}&=A_{9,*}+\frac{7\lambda}{M}
             -\frac{2863}{240M^2}+\frac{32}{M^3},
 \label{s5-final-A}\\
 \limsup_{\substack{K\to\infty\\40\mid K}}
 K^{-2}\log m_{9,K,M}&\leq A_{9,M}\qquad(40\mid M).
 \notag
\end{align}
The last inequality is conditional on the local estimates, as before.
The coefficient $7\lambda/M$ survives because the new small-prime cost
$10\lambda/M$ is offset by the tail's $3\lambda/M$.
An especially transparent comparison with the source is
\begin{equation}
 A_{9,M}-A_{5,M}
 =\frac{45541}{57600}
    +4\lambda\left(\frac13-\frac1M\right)+\frac{4\lambda}{M}
 =\frac{116581}{57600}.
 \label{s5-budget-increment}
\end{equation}
All displayed rational integrals are evaluated algebraically.  They
describe the size of this proposed normalization; they do not supply a
missing decay inequality or assert irrationality.

\section{The real determinant estimate}
\label{s6-real}

We follow source Section~6, pp.~17--20, retaining $\alpha=3/40$,
$\lambda=37/40$, $D_N^6$, and $S_K$.  Consequently the leading external
field is literally the same function:
\begin{equation}
 V(t)=2\pi\sqrt t+\int_0^1\log(t+u^2)\,du
             -6\int_0^\alpha\log(t+u^2)\,du.
 \label{s6-field}
\end{equation}
Write $U_\rho(t)=\int\log|t-u|\,d\rho(u)$ and
$I(\rho)=\iint\log|t-u|\,d\rho(t)d\rho(u)$.  The source's rational
sixteen-component measure has mass $\lambda$, support in $(0,2)$, and
component lengths exceeding $1/225$.  Its numerical certificate is the
following \textbf{[Formal analogue: reused source input]}, corresponding
to (6.2)--(6.4):
\begin{align}
 2U_\rho(t)-V(t)&\le M_0=-\frac{1329}{200},\quad t\ge0,
 \label{s6-potential}\\
 C_*&=-2\lambda+12\alpha\lambda(1-\log\alpha)
          +3\lambda^2-2\lambda^2\log(2\lambda),\\
 \lambda M_0-I(\rho)+C_*&\le U=-\frac{2733991}{2000000}.
 \label{s6-source-constant}
\end{align}
Appendix~\ref{appa-measure} gives the parameters and the certificate's
procedure.  Reusing it is justified by equality of $V$; this draft does
not independently rerun the source's interval checks.

\subsection{A bound for every configuration}

\textbf{[Derived]} The zero-mass energy argument of source Lemma~6.2
does not involve the zeta index.  If a compactly supported real signed
measure $\nu$ has mass zero and
$\iint|\log|z-w||\,d|\nu|(z)d|\nu|(w)<\infty$, completing the square gives
\[
 e^{-s|z-w|^2}=\frac{4s}{\pi}\int_{\mathbb C}
 e^{-2s|z-u|^2}e^{-2s|w-u|^2}\,dA(u),
\]
and hence
\[
 J(s):=\iint e^{-s|z-w|^2}\,d\nu(z)d\nu(w)
 =\frac{4s}{\pi}\int_{\mathbb C}
       \left(\int e^{-2s|z-u|^2}\,d\nu(z)\right)^2dA(u)\ge0.
\]
For $0<a<b$, put
$L_{a,b}(r)=\tfrac12\int_a^b(e^{-s}-e^{-sr^2})\,ds/s$.
The full integral is $\log r$, its integrand has constant sign, and
$|L_{a,b}(r)|\le|\log r|$.  Mass zero therefore implies
\[
 \iint L_{a,b}(|z-w|)\,d\nu(z)d\nu(w)
       =-\frac12\int_a^b\frac{J(s)}s\,ds\le0.
\]
The integrability hypothesis removes the diagonal and permits dominated
convergence.  Thus $I(\nu)\le0$, the analogue of (6.5).

For $t_1,\ldots,t_h\ge0$, let $\omega_i$ be unit arclength measure on
the complex circle of radius $\varepsilon$ about $t_i$, and set
$\sigma=K^{-1}\sum_i\omega_i$.  Then $\sigma-\rho$ has mass zero.
Circle and arcsine potentials are continuous and finite; boundedness of
the positive logarithmic kernel on their joint compact support then gives
the required absolute integrability.  The circle-average formula is
\[
 \int\log|z-u|\,d\omega_i(z)=\log\max(|t_i-u|,\varepsilon).
\]
A unit arcsine measure on an interval of length $L$ assigns at most
$4\sqrt{2r}/(\pi\sqrt L)$ to an interval of radius $r$.  Integrating
this bound with $dr/r$ gives
\[
 \int U_\rho\,d\omega_i-U_\rho(t_i)
 \le\sum_{j=1}^{16}\frac{8c_j\sqrt{2\varepsilon}}
                              {\pi\sqrt{b_j-a_j}}
 \le60\sqrt\varepsilon.
\]
Moreover $I(\omega_i)=\log\varepsilon$, and the mutual circle energy
is at least $\log|t_i-t_j|$.  Expanding $I(\sigma-\rho)\le0$ and
using \eqref{s6-potential} proves the analogue of (6.6):
\begin{equation}
 2\sum_{i<j}\log|t_i-t_j|-K\sum_iV(t_i)
 \le(\lambda M_0-I(\rho))K^2
       +120\lambda K^2\sqrt\varepsilon-h\log\varepsilon.
 \label{s6-smoothed}
\end{equation}
Coincident points give a left side of $-\infty$.

To integrate over an unbounded domain we must retain decay.  For $t\ge2$,
support of $\rho$ and $\log(1+x)\le x$ give
\begin{equation}
 2U_\rho(t)-V(t)
 \le(2\lambda-1+6\alpha)\log t+\frac{2\alpha^3}{t}-2\pi\sqrt t
 \le\frac{13}{10}\log t+\frac{6\alpha^3}{t}-2\pi\sqrt t.
 \label{s6-tail}
\end{equation}
After adding $\sqrt t/K$, the last expression decreases for $t\ge2$,
$K\ge2$, and is below $M_0$ at $2$.  For $t\le2$ use
\eqref{s6-potential}.  Thus
$2U_\rho-V+\sqrt t/K\le M_0+\sqrt2/K$, as in (6.7)--(6.8).
Apply \eqref{s6-smoothed} with this modified field and
$\varepsilon=K^{-2}$ to obtain
\begin{equation}
 2\sum_{i<j}\log|t_i-t_j|-K\sum_iV(t_i)+\sum_i\sqrt{t_i}
 \le(\lambda M_0-I(\rho))K^2+(120+\sqrt2)h+2h\log K.
 \label{s6-configuration}
\end{equation}

\subsection{The Gram integral and scaling}

\textbf{[Derived]} Expanding the two Vandermonde determinants and
integrating term by term gives Andr\'eief's identity, now with $w_9$:
\begin{equation}
 \Delta_{9,K}(\zeta(9))=\frac1{h!}\int_{(0,\infty)^h}
 \prod_{i<j}(y_i^2-y_j^2)^2
 \prod_{i=1}^h\frac{D_N(y_i^2)^6}{D_K(y_i^2)}w_9(y_i)\,dy_i.
 \label{s6-andreief}
\end{equation}
This is the analogue of (6.10); the positive weight gives strict
positivity, and exponential decay gives convergence.  With $q=e^{-2\pi y}$,
\[
 \sum_{\ell\ge1}\ell^8q^\ell
 =\frac{q(1+247q+4293q^2+15619q^3+15619q^4
                         +4293q^5+247q^6+q^7)}{(1-q)^9}.
\]
The numerator's coefficients sum to $40320$.  Using
$1-e^{-u}\ge u/(1+u)$ and $\pi<4$ therefore proves the replacement
for (6.11):
\begin{equation}
 w_9(y)\le\frac{(1+2\pi y)^9}{\pi}e^{-2\pi y}
          \le C_9(1+y)^9e^{-2\pi y},\qquad C_9=2^{25}.
 \label{s6-weight-bound}
\end{equation}

For $1\le m\le K$, monotonicity, followed by the factorial bound at
$t=0$, yields the unchanged Riemann-sum error (6.12):
\begin{equation}
 0\le\sum_{j=1}^m\log(t+(j/K)^2)
       -K\int_0^{m/K}\log(t+u^2)\,du\le2\log K+2.
 \label{s6-riemann}
\end{equation}
In particular the six numerator sums contribute at most $12\log K+12$;
the denominator error has the favorable sign.  Put $y_i=K\sqrt{t_i}$
in \eqref{s6-andreief}.  The powers from the Vandermonde, rational
factor, weight bound, and Jacobian are respectively
$2h(h-1)$, $(12N-2K)h$, $9h$, and $h$.  Including the Riemann-sum
loss $12h$ gives the analogue of (6.13):
\begin{align}
 \Delta_{9,K}(\zeta(9))&\le
 \frac{K^{2h(h+6N-K)+20h}(C_9e^{12}/2)^h}{h!}\notag\\
 &\quad\times\int_{(0,\infty)^h}\prod_{i<j}(t_i-t_j)^2
       \prod_i e^{-KV(t_i)}t_i^{-1/2}(1+\sqrt{t_i})^9\,dt_i.
 \label{s6-scaled}
\end{align}
The remaining one-variable integral is explicit:
\[
 R_9=\int_0^\infty t^{-1/2}(1+\sqrt t)^9e^{-\sqrt t}\,dt
     =2\int_0^\infty(1+u)^9e^{-u}\,du
     =2\sum_{k=0}^9\binom9k k!=1972820.
\]
Apply \eqref{s6-configuration}; since
$\log(C_9/2)+12+\log R_9+120+\sqrt2<170$, dropping $-\log h!$
gives
\begin{equation}
 \log\Delta_{9,K}(\zeta(9))\le
 2h(h+6N-K)\log K+(\lambda M_0-I(\rho))K^2
                 +22h\log K+170h.
 \label{s6-delta-bound}
\end{equation}

The normalization also remains algebraically identical.  Expanding
\[
 \log S_K=2h\log(K!)-12h\log(N!)+(h-1)\log4
                 -2\sum_{i=1}^{h-1}\log((2i)!)
\]
with $m\log m-m+1\le\log(m!)\le m\log m-m+\log m+2$ and
$-2\sum_{i<h}\log((2i)!)\le-2h^2\log(2h)+3h^2+4h\log(2h)$
gives source (6.15):
\[
 \log S_K\le(2\lambda-12\alpha\lambda-2\lambda^2)K^2\log K
                         +C_*K^2+6h\log K+6h.
\]
The quadratic logarithmic terms cancel because
$2\lambda(\lambda+6\alpha-1)+2\lambda-12\alpha\lambda-2\lambda^2=0$.
Thus the real-estimate conclusion is
\begin{equation}
 \boxed{\quad 0<F_{9,K}(\zeta(9)),\qquad
 \log F_{9,K}(\zeta(9))\le UK^2+28K\log K+200K.\quad}
 \label{s6-final}
\end{equation}
\textbf{[Conditional]} The displayed numerical $U$ uses the source
Appendix~A certificate as input.  All changes caused by $w_9$ are the
explicit weight and scaling computations above.  This real decay alone
does not assert decay after arithmetic denominator clearing.

\section{Combining the arithmetic and real estimates}
\label{s7-completion}

This section is the analogue of source Section~7, p.~21.  Its task is to
put the two quadratic coefficients on the same scale.  We retain
$K=40n$, $N=3n$, $h=37n$, and write $\xi=\zeta(9)$.  The outcome is a
positive upper-bound budget; the final small-value step of the source
therefore becomes a conditional implication.

\subsection{The normalized polynomial and its budget}
\label{s7-budget}

\noindent\textbf{[Conditional]} With the arithmetic estimates used in
Section~5, let
\begin{equation}
 Q_{9,K,M}(X)=m_{9,K,M}F_{9,K}(X)
            =m_{9,K,M}S_K\Delta_{9,K}(X)\in\mathbb Z[X],
 \qquad m_{9,K,M}>0.
 \label{s7-normalized}
\end{equation}
The degree is $h$ by the leading-coefficient calculation in Section~2.
The positive Gram representation and the positive scalar factors give
$Q_{9,K,M}(\xi)>0$.  The normalization changes neither fact.  For fixed
$M\in40\mathbb Z_{>0}$ the input estimates are
\[
 \log m_{9,K,M}\leq A_{9,M}K^2+o_M(K^2),\qquad
 \log F_{9,K}(\xi)\leq UK^2+28K\log K+200K,
 \qquad U=-\frac{2733991}{2000000}.
\]
The second estimate is \eqref{s6-final}; it uses the source Appendix~A
numerical certificate discussed in Appendix~\ref{appa-measure}.
Adding, and then dividing by $K^2$, gives the
direct analogue of source (7.1):
\begin{equation}
 \limsup_{\substack{K\to\infty\\40\mid K}}
  \frac{\log Q_{9,K,M}(\xi)}{K^2}
 \leq B_{9,M},\qquad B_{9,M}:=A_{9,M}+U.
 \label{s7-combined}
\end{equation}
Thus the arithmetic loss has to be paid before the negative real
coefficient can be used to obtain a small integer polynomial value.
The order of the limits matters here: $M$ is held fixed while $K$
increases through multiples of forty.  The comparisons between cutoffs
below compare the resulting constants.  They do not require a uniform
error estimate with a cutoff $M$ depending on $K$.  In particular, the
explicit $K\log K$ error in the real estimate does not upgrade the
arithmetic remainder $o_M(K^2)$ to an error of that same order.

\noindent\textbf{[Derived]} The arithmetic coefficient obtained with the
unchanged $D_N^6$ construction is
\begin{align}
 A_{9,M}&=A_{5,M}+\frac{116581}{57600}
 \notag\\
 &=A_{9,*}+\frac{7\lambda}{M}
       -\frac{2923/240-1/4}{M^2}+\frac{32}{M^3},
 \qquad \lambda=\frac{37}{40},
 \label{s7-arithmetic-coefficient}\\
 A_{9,*}&=\frac{25180312106351244019069}
                   {7535670527041937280000}
           \simeq3.341482621353.
 \label{s7-limit-coefficient}
\end{align}
Here $A_{5,M}$ denotes the source coefficient, rather than a new
fifth-order construction.  The rational shift summarizes the increased
denominator cost calculated in Section~5.  In particular, the comparison
replacing the negative margins of source (7.2) is
\begin{equation}
 \begin{array}{c|c|c}
 M&A_{9,M}&B_{9,M}=A_{9,M}+U\\ \hline
 200&3.373563392186\ldots&2.006567892186\ldots\\
 100000&3.341547370160\ldots&1.974551870160\ldots\\
 \infty&3.341482621353\ldots&1.974487121353\ldots
 \end{array}
 \label{s7-sign-table}
\end{equation}
Increasing the cutoff does not change this sign.  Indeed, the correction
to $A_{9,*}$ is
\[
 \frac{1554M^2-2863M+7680}{240M^3}>0\qquad(M\geq40),
\]
and even $A_{9,*}+U$ is positive.

\subsection{The conditional final step}
\label{s7-conditional-step}

\noindent\textbf{[Blocked conclusion]} Substitution of $K=40n$ in
\eqref{s7-combined} supplies an upper coefficient $1600B_{9,M}>0$ on
the $n^2$ scale.  It supplies no decay rate of the type used in source
(2.7)--(2.8).  This does not prove that the polynomial values grow: a
positive upper bound alone does not decide their actual size.  It says
precisely that the present arithmetic and real estimates do not close
the small-value argument.

\noindent\textbf{[Conditional]} The remaining manipulation can still be
recorded.  Suppose some strengthened estimates, for a fixed cutoff,
gave an $\varepsilon>0$ and
\begin{equation}
 0<Q_{9,K,M}(\xi)\leq e^{-\varepsilon K^2}
 \quad\text{for all sufficiently large }K\in40\mathbb Z_{>0}.
 \label{s7-hypothetical-decay}
\end{equation}
If $\xi=a/b$ with $a\in\mathbb Z$ and $b\in\mathbb Z_{>0}$, degree
$h=\lambda K$ and integrality would imply
\[
 0<b^hQ_{9,K,M}(a/b)
 \leq\exp\{\lambda K\log b-\varepsilon K^2\}\longrightarrow0,
 \qquad b^hQ_{9,K,M}(a/b)\in\mathbb Z.
\]
This is the source's final contradiction, with the moment order changed
to nine.  Its algebra is unchanged, but its negative-budget hypothesis
is not furnished by \eqref{s7-sign-table}.  Appendix~C likewise records
the rational-approximation implication under an explicit
$\varepsilon$ hypothesis, while retaining the relative norm and height
calculations that can be made with the present budget.

\appendix
\section{The comparison measure and its elementary bounds}
\label{appa-measure}

This is the analogue of source Appendix~A, pp.~21--26.  The measure and
the numerical enclosures are \textbf{[Formal analogue: reused source
input]}; their defining algebra is displayed below.  No new search for
an equilibrium measure or independent verification of the source's
684 partition intervals is claimed.

\subsection{The sixteen intervals}

Let $d\omega_{[a,b]}(t)=dt/(\pi\sqrt{(t-a)(b-t)})$ on $(a,b)$ and set
$\rho=\sum_{j=1}^{16}c_j\omega_{[a_j,b_j]}$.  The following integers,
each divided by $10^{12}$, reproduce source Table~1:
\begin{center}
\begin{tabular}{r|r|r|r}
$j$&$10^{12}a_j$&$10^{12}b_j$&$10^{12}c_j$\\\hline
1&3906748086&8992695531&10515596180\\
2&2312248264&15340997855&29471737793\\
3&1402286665&25730180724&42934365099\\
4&881725356&41909578246&58204231966\\
5&578197906&65851089563&69037621310\\
6&396324613&99481037884&78873099189\\
7&283911191&144325727458&84856120711\\
8&212206188&201105762729&88396082127\\
9&165097686&269345996903&88303382125\\
10&133347132&347089554156&85472321255\\
11&111522114&430806704415&78899184238\\
12&96349355&515561896511&70353471918\\
13&85815639&595448778546&58838976615\\
14&78667711&664241383483&44421321106\\
15&74129565&716160577112&30462865791\\
16&71741310&746637295669&5959622577
\end{tabular}
\end{center}
Direct rational arithmetic gives $\sum c_j=37/40$,
$0<a_{16}<\cdots<a_1<b_1<\cdots<b_{16}<2$, and
$b_j-a_j>1/225$.  With $m=(a+b)/2$, the substitution
$u=m+(b-a)\cos\theta/2$ gives
\begin{equation}
 U_{\omega_{[a,b]}}(t)=
 \begin{cases}
 \log((b-a)/4),&a\le t\le b,\\
 \displaystyle\log\frac{|t-m|+\sqrt{(t-a)(t-b)}}2,&t\notin[a,b].
 \end{cases}
 \label{appa-arcsine}
\end{equation}
For the exterior branch, differentiate the defining integral and
normalize by $U(t)-\log|t|\to0$; on the interval, the sine-product
identity for $\cos\phi-\cos\theta$ gives the constant.  This proves
the algebra of (A.1).

Put $S_j=\sum_{i\le j}c_i$, $S_0=0$.  An outer interval's potential
is constant on every inner interval, so the self and cross terms combine
as $c_j^2+2c_jS_{j-1}=S_j^2-S_{j-1}^2$.  Hence (A.2) becomes
\begin{equation}
 I(\rho)=\sum_{j=1}^{16}(S_j^2-S_{j-1}^2)
                          \log\frac{b_j-a_j}{4}.
 \label{appa-energy}
\end{equation}

\subsection{Elementary enclosures}

For $1\le x\le2$, $z=(x-1)/(x+1)$, the logarithmic series gives (A.3):
\[
 0\le\log x-2\sum_{k=0}^{m-1}\frac{z^{2k+1}}{2k+1}
       \le\frac{2z^{2m+1}}{(2m+1)(1-z^2)}.
\]
Reduce a general argument as $x=2^er$, $1\le r<2$.  For
$0\le z\le1/2$, the alternating series gives (A.4):
\[
 \left|\arctan z-\sum_{k=0}^{m-1}
                  \frac{(-1)^kz^{2k+1}}{2k+1}\right|
       \le\frac{z^{2m+1}}{2m+1}.
\]
Use $\arctan x=\pi/2-\arctan(1/x)$ for $x>1$ and
$\arctan x=2\arctan(x/(1+\sqrt{1+x^2}))$ for $0\le x\le1$;
enclose $\pi$ through $16\arctan(1/5)-4\arctan(1/239)$.
The source chooses $m=64$ for logarithms, $m=80$ for arctangents,
and adjacent dyadic multiples of $2^{-144}$ for square roots.
The reduced series errors are below $2^{-144}$; interval arithmetic
must also propagate the rounding errors through each expression.

\subsection{The potential inequality on the half-line}

Integrating the field in \eqref{s6-field} gives (A.5):
\begin{align}
 V(t)&=\log(1+t)-6\alpha\log(t+\alpha^2)-2+12\alpha\notag\\
 &\quad+2\sqrt t\left(\pi+\arctan\frac1{\sqrt t}
                                -6\arctan\frac\alpha{\sqrt t}\right),
 \label{appa-field-explicit}
\end{align}
with $V(0)=-12\alpha\log\alpha-2+12\alpha$.
For $\Phi(y)=V(y^2)$,
\[
 \Phi'(y)=2\left(\pi+\arctan(1/y)-6\arctan(\alpha/y)\right),
 \qquad \Phi''(y)=-\frac2{1+y^2}+\frac{12\alpha}{\alpha^2+y^2}.
\]
The latter changes sign at $y^2=711/880$, while
$\Phi'(0+)=-3\pi$ and $\Phi'(\infty)=2\pi$.  Thus $V$ has one minimum.
The source encloses it between
$q_-=59205077/10^{10}$ and $q_+=59205079/10^{10}$ by opposite
signs of $\Phi'(\sqrt{q_\pm})$.
On this small interval a lower bound, as in (A.6), is
\begin{align*}
 V_*={}&\log(1+q_-)-6\alpha\log(q_++\alpha^2)-2+12\alpha\\
 &+2\sqrt{q_+}\left(\pi+\arctan\frac1{\sqrt{q_+}}
                              -6\arctan\frac\alpha{\sqrt{q_-}}\right).
\end{align*}
The parenthesis is negative, which determines the direction of the
square-root bound.  Nesting makes $U_\rho$ nonincreasing up to $a_1$,
constant on $[a_1,b_1]$, and nondecreasing thereafter.  Thus (A.7) is
\[
 B(l,r)=2\max\{U_\rho(l),U_\rho(r)\}
 -\begin{cases}V(r),&r\le q_-,\\V(l),&l\ge q_+,\\V_*,&\text{otherwise},\end{cases}
 \qquad 2U_\rho(t)-V(t)\le B(l,r)\quad(l\le t\le r).
\]
For completeness, the partition rule (A.8) is
\[
 A_0=0,\quad A_j=a_{17-j}\ (1\le j\le16),\quad
 A_{17}=q_-,\quad A_{18}=q_+,\quad
 A_{18+j}=b_j\ (1\le j\le16),\quad A_{35}=2,
\]
where a source Table~2 row $(j,d,k)$ specifies
\[
 \left[A_j+(A_{j+1}-A_j)\frac{k}{2^d},\quad
       A_j+(A_{j+1}-A_j)\frac{k+1}{2^d}\right].
\]
Use exactly the row list in source Table~2, pp.~24--25.  Its claimed
rational certificate (A.9) is
\[
 B(l,r)<-\frac{6645002}{10^6}<-\frac{1329}{200}
\]
on every listed interval.  The tail estimate \eqref{s6-tail} treats
$t\ge2$.  Because both $\rho$ and $V$ are unchanged, these are the same
numerical inputs in the $\zeta(9)$ draft, not new checks.

\subsection{Energy and the norm constant}

Applying the stated logarithmic enclosures to \eqref{appa-energy} and
$C_*$ gives the source input (A.10):
\[
 -\frac{2126593445148}{10^{12}}<I(\rho)
              <-\frac{2126593445147}{10^{12}},\qquad
 \frac{2653035990340}{10^{12}}<C_*
              <\frac{2653035990341}{10^{12}}.
\]
Since $\lambda M_0=-49173/8000$, the directions of the bounds give
\[
 \lambda M_0-I(\rho)+C_*<-
 \frac{1366995564511}{10^{12}}<-
 \frac{2733991}{2000000}=U.
\]
Finally, (A.11) encloses the source's norm exponent
\[
 \kappa=\frac{2\pi}{\sqrt3}+\log\frac43+\frac9{20}\log18
 +\frac{37}{40}\left(1+\log\frac{1600}{111}\right)
 +\frac{37}{20}\log\frac{43+4\sqrt{114}}5
\]
by
\[
 \frac{2773171335773}{200000000000}<\kappa
 <\frac{6932928339433}{500000000000}<\frac{6933}{500}.
\]
These rational enclosures remain source input.  The reason the same
leading norm exponent is available for the fixed-shape analogue is that
the lower bound from the $\ell=1$ term of $w_9$ changes powers and
constants, whereas its exponential factor remains $e^{-2\pi y}$;
the separate norm manipulation belongs to Appendix~C.

\section{Exact arithmetic integrals}
\label{appb-integrals}

This appendix follows source Appendix~B, pp.~26--28.  Its calculations
are \textbf{[Derived]} identities for the functions in Section~5;
their use as determinant bounds has the same local dependencies.

\subsection{Evaluation of the inner integral}
\label{appb-inner}

For $0\leq u<1$, set $d_0(u)=\min(u,1-u)$ and let $e(u)$ equal
$1$ for $u\leq1/2$ and $-1$ otherwise.  For $f=\{x\}$,
\[
 \ell(x,z)-2x=e(f)\bigl(\boldsymbol{1}_{z<d_0(f)}-2d_0(f)\bigr)
\]
away from a single irrelevant endpoint.  Hence source (B.1) remains
exactly applicable:
\begin{align}
 \int_0^{1/2}B^2\,dz&=d_0(g)(1-2d_0(g)),\notag\\
 \int_0^{1/2}AB\,dz&=e(f)e(g)
       \bigl(\min(d_0(f),d_0(g))-2d_0(f)d_0(g)\bigr).
 \label{appb-fractional-parts}
\end{align}
Together with \eqref{s5-tail-decomposition}, this is an explicit
rational formula for $R_9$ without a remaining $z$-integral.

Let $3=t_0<\cdots<t_{143}=20$ be the endpoints and the interior
points where $cx$ is integral for some
\[
 c\in\{2,2\alpha,2\lambda,2H,4\alpha,2(1-\alpha),2(1+\alpha)\}.
\]
These are the same breakpoints as source (B.2).  The changes of branch
in the minima and positive parts are included.  The quadratic
coefficient still vanishes because
\[
 2\lambda-12\lambda\alpha
 +2(H^2-H-\lambda^2)-18\alpha^2+6\alpha=0.
\]
Thus $R_9=a_i x+b_i$ on each open interval.  Taking its values at
$x_i=(2t_i+t_{i+1})/3$ and $y_i=(t_i+2t_{i+1})/3$ determines
$a_i=(R_9(y_i)-R_9(x_i))/(y_i-x_i)$ and
$b_i=R_9(x_i)-a_i x_i$.  The exact analogue of (B.3) is
\begin{equation}
 \int_3^{20}\frac{R_9(x)}{x^3}\,dx
 =\sum_{i=0}^{142}\left[
 a_i\left(\frac1{t_i}-\frac1{t_{i+1}}\right)
 +\frac{b_i}{2}\left(\frac1{t_i^2}-\frac1{t_{i+1}^2}\right)
 \right].
 \label{appb-inner-sum}
\end{equation}
Exact rational evaluation of the original floor-function expression
for $\Gamma_9$ gives \eqref{s5-finite-inner}; a third interior point
on each interval checks the affine interpolation.  There is also a
short independent comparison: $a_{9,i}=a_{5,i}+4\lambda$ and
$b_{9,i}=b_{5,i}$, so the increase is
\[
 4\lambda\left(\frac13-\frac1{20}\right)=\frac{629}{600}.
\]
In particular, every entry of source Table~3 has the explicit analogue
\[
 \int_j^{j+1}R_9(x)x^{-3}\,dx
 =\int_j^{j+1}R_5(x)x^{-3}\,dx+\frac{37}{10j(j+1)},
 \qquad j=3,\ldots,19.
\]
This gives all seventeen unit-interval values from one exact rule.

\subsection{Evaluation of the outer integral}
\label{appb-outer}

Write $T_9(y)=C_9(y)+S(y)$.  Resolving the minima, positive parts,
and floor in Section~5 gives the following replacement for source
Table~4; on each interval $T_9(y)=b+cy$.
\begin{center}
\begin{tabular}{rrrr}
\hline
$l$&$r$&$b$&$c$\\
\hline
$1/3$&$43/120$&$35/4$&$-8$\\
$43/120$&$37/100$&$76/5$&$-26$\\
$37/100$&$37/80$&$267/20$&$-21$\\
$37/80$&$1/2$&$119/20$&$-5$\\
$1/2$&$43/80$&$39/5$&$-5$\\
$43/80$&$37/60$&$57/4$&$-17$\\
$37/60$&$37/40$&$62/5$&$-14$\\
$37/40$&$1$&$5$&$-6$\\
$1$&$37/20$&$37/20$&$-1$\\
\hline
\end{tabular}
\end{center}
The old rank-related breakpoints $13/30$ and $13/20$ disappear.
Summing $b(r-l)+c(r^2-l^2)/2$ gives $305479/144000$.
Separately, the raw determinant integrals satisfy
\[
 \int C_9=\frac{14441}{7200},\qquad
 \int C_5=\frac{23329}{19200},\qquad
 \int(C_9-C_5)=\frac{45541}{57600},
\]
where all three integrals run from $1/3$ to $1$ and
$C_5=R_0-d$ is the source's cost.  The common scalar integral
$1851/16000$ provides an additional check on the comparison.

\subsection{The final rational margins}
\label{appb-margins}

\paragraph{[Blocked conclusion] The source's margins change sign.}
With the unchanged leading real coefficient
$U=-2733991/2000000$, the exact sums at the two source cutoffs are
\begin{align*}
 A_{9,200}+U
 &=\frac{189010431570697800693281}{94195881588024216000000}
   =2.006567892\ldots,\\
 A_{9,100000}+U
 &=\frac{11624665884437287177113589319057}
          {5887242599251513500000000000000}
   =1.974551870\ldots.
\end{align*}
Consequently the proposed decay coefficients $-1600(A_{9,M}+U)$
are negative.  In the exact format of source Appendix~B.3, the
differences from its target margins are
\begin{align*}
 -1600(A_{9,200}+U)-\frac{139}{5}
 &=-\frac{190647085013289721446281}{58872425992515135000},\\
 -1600(A_{9,100000}+U)-\frac{7907}{100}
 &=-\frac{11915606054639047909891714319057}
           {3679526624532195937500000000}.
\end{align*}
The limiting sum is also positive:
$A_{9,*}+U=1.974487121\ldots$.
These are failures of the required sign for this certificate.
A positive upper bound is not a lower bound on the actual polynomial
values, so none of these calculations asserts growth, rationality,
or impossibility of another construction.

\section{Rational approximation and coefficient height}
\label{c-approximation}

This appendix follows source Appendix~C, pp.~28--31, including all steps
of its (C.1)--(C.8).  The matrix perturbation calculation survives the
change of moment order.  We first derive its relative norm estimate,
then show exactly where a negative completion budget would be required.
Throughout, $K=40n$, $N=3n$, $h=37n$, $\alpha=3/40$, and $\lambda=37/40$.

\subsection{The relative matrix norm}
\label{c-relative-norm}

Write $\xi=\zeta(9)$, $G=G_{9,K}(\xi)$, and
$M_{9,K}=[X]G_{9,K}(X)$.  The positive Gram representation makes $G$
positive definite.  Consequently the real symmetric matrix
$B_K=G^{-1/2}M_{9,K}G^{-1/2}$ is defined.  Its Euclidean operator norm
measures a change of $X$ relative to the size of the unperturbed matrix.
The analogue of source (C.1) is
\begin{equation}
 \|B_K\|\leq e^{\rho K},\qquad
 \rho=\frac{6933}{500}=13.866,
 \quad K\in40\mathbb Z_{>0}\text{ sufficiently large}.
 \label{c-norm-bound}
\end{equation}
The value of $\rho$ is the source's fraction $6933/500$; its decimal is
$13.866$.

\noindent\textbf{[Derived]} Let $q$ be a nonzero real polynomial of
degree less than $h$, write $\mathbf q$ for its coefficient vector, and
put $p(t)=q(K^2t)$.  For fixed $0<a<b$, retain just the $\ell=1$ term in
the positive weight
\[
 w_9(y)=\frac{(2\pi)^8y^9}{20160}
          \sum_{\ell\geq1}\ell^8e^{-2\pi\ell y}.
\]
With $y=K\sqrt t$, this term contributes
$K^{10}(2\pi)^8t^4e^{-2\pi K\sqrt t}\,dt/40320$.
On $[a,b]$ the scaled numerator is at least $a^{6N}$ and the scaled
denominator at most $(1+b)^K$.  The analogue of source (C.2) is therefore
\begin{equation}
 \mathbf q^TG\mathbf q\geq
 A_Kc_a e^{-F(a,b)K}\int_a^b p(t)^2\,dt,
 \quad
 \begin{aligned}
 A_K&=K^{12N-2K+10},&
 c_a&=\frac{(2\pi)^8a^4}{40320},\\
 F(a,b)&=2\pi\sqrt b+\log(1+b)-6\alpha\log a.
 \end{aligned}
 \label{c-gram-lower}
\end{equation}
The extra four powers of $K$ relative to the source occur in both
quadratic forms below, so they cancel from their ratio.
Restricting the positive integral is useful even though it discards
most of the measure: it controls the entire coefficient vector through
the squared norm of $p$ on one fixed interval.  No cancellation enters
this lower bound.  The later Legendre estimate converts that interval
norm into control of the evaluations at the negative pole locations.

Only pole terms contribute to $M_{9,K}$, since polynomial moments do
not depend on $X$.  Cancel one factor $D_N$ from $D_N^6/D_K$ and take
the residues at $-j^2$, $N<j\leq K$.  With $u_j=j/K$, source (C.3)
becomes
\begin{equation}
 \begin{aligned}
 \mathbf q^TM_{9,K}\mathbf q
    &=A_K\sum_{j=N+1}^{K}r_jp(-u_j^2)^2,\\
 r_j&=\frac{u_j^8\displaystyle\prod_{\ell=1}^{N}
                      (u_\ell^2-u_j^2)^5}
                 {\displaystyle\prod_{\substack{N<k\leq K\\k\ne j}}
                      (u_k^2-u_j^2)}.
 \end{aligned}
 \label{c-residues}
\end{equation}
For clarity, before scaling the residue multiplier is
$j^8D_N(-j^2)^5/\prod_{N<k\leq K,\,k\ne j}(k^2-j^2)$.
Its power of $K$ is
$8+10N-2(h-1)=12N-2K+10$, confirming the same $A_K$ as in the integral.
The factor $u_j^8$ is the sole change in this scaled residue formula
from the fourth power appearing in the source.
The numerator has absolute value at most one.  For each denominator
factor, $u_k+u_j\geq2\alpha$ and
$|u_k-u_j|=|k-j|/K$.  Setting $m=h-1$ gives the analogue of source (C.4):
\begin{equation}
 \begin{aligned}
 |r_j|&\leq
 \frac{(K/(2\alpha))^{h-1}}{(j-N-1)!(K-j)!}
 \leq\frac{(K/\alpha)^{h-1}}{(h-1)!}\\
 &\leq\exp\{\lambda K(1-\log(\alpha\lambda))\}.
 \end{aligned}
 \label{c-residue-bound}
\end{equation}
The middle step uses $\binom m{j-N-1}\leq2^m$.  For the last step,
$m!\geq(m/e)^m$ and the function $x\log(eK/(\alpha x))$ is increasing
on $0<x\leq h$.

Expand $p$ in the Legendre polynomials for $[a,b]$.  At $-1\leq t\leq0$
their transformed argument has absolute value at most
$Y=(2+a+b)/(b-a)$.  The source's Legendre bound $|P_j(s)|\leq R^j$,
where $R=Y+\sqrt{Y^2-1}$, and Cauchy--Schwarz give
\begin{equation}
 \max_{-1\leq t\leq0}|p(t)|^2\leq
 \frac{h^2R^{2h}}{b-a}\int_a^b p(t)^2\,dt,
 \qquad Y=\frac{2+a+b}{b-a},\quad R=Y+\sqrt{Y^2-1}.
 \label{c-legendre}
\end{equation}
Indeed, the squared norms are $(b-a)/(2j+1)$ and
$\sum_{j=0}^{h-1}(2j+1)=h^2$.  This is exactly the step in source (C.5).
There are $h$ residues, so \eqref{c-gram-lower}--\eqref{c-legendre} imply
\begin{equation}
 \begin{aligned}
 \|B_K\|&=
  \sup_{\mathbf q\ne0}
  \frac{|\mathbf q^TM_{9,K}\mathbf q|}{\mathbf q^TG\mathbf q}
  \leq\frac{h^3}{c_a(b-a)}e^{\kappa(a,b)K},\\
 \kappa(a,b)&=F(a,b)+\lambda(1-\log(\alpha\lambda))+2\lambda\log R.
 \end{aligned}
 \label{c-kappa-bound}
\end{equation}
For the source choice $a=1/18$, $b=1/3$, we have $Y=43/5$ and
$R=(43+4\sqrt{114})/5$.  Thus source (C.6) is unchanged:
\begin{equation}
 \begin{aligned}
 \kappa={}&\frac{2\pi}{\sqrt3}+\log\frac43
       +\frac9{20}\log18
       +\frac{37}{40}\left(1+\log\frac{1600}{111}\right)
       +\frac{37}{20}\log\frac{43+4\sqrt{114}}5\\
   &\simeq13.865856678865<\frac{6933}{500}.
 \end{aligned}
 \label{c-kappa-value}
\end{equation}
\noindent\textbf{[Conditional]} Using the elementary enclosure from
source Appendix~A for this strict numerical inequality, its positive
margin absorbs $h^3/(c_a(b-a))$ for all sufficiently large $K$, proving
\eqref{c-norm-bound}.  The exponent is unchanged; the threshold need
not be the source's threshold.
More precisely, the retained enclosure is
$\kappa<6932928339433/500000000000<6933/500$ from source (A.11).
Absorption uses only $\log(h^3/(c_a(b-a)))=O(\log K)$, whereas the
positive gap between the two exponents is multiplied by $K$.

\subsection{Perturbation and the conditional approximation statement}
\label{c-approximation-conditional}

Fix a cutoff $M$ and abbreviate $Q_K=Q_{9,K,M}$.  Affine dependence
$G_{9,K}(\xi+\delta)=G+\delta M_{9,K}$ gives the exact identity
\begin{equation}
 Q_K(\xi+\delta)=Q_K(\xi)\det(I+\delta B_K)
      =Q_K(\xi)\prod_{j=1}^{h}(1+\delta\beta_j),
 \label{c-determinant-perturbation}
\end{equation}
where $\beta_j$ are the real eigenvalues of $B_K$.  If
$|\delta|\|B_K\|<1$, every factor is positive.  Moreover
\[
 \det(I+\delta B_K)
 \leq(1+|\delta|\|B_K\|)^h
 \leq\exp\{h|\delta|\|B_K\|\}.
\]
This explains how a polynomial that is small at $\xi$ can be kept
nonzero at a sufficiently close rational number.
The scalar $m_{9,K,M}S_K$ cancels in this relative comparison; it is
independent of $X$.  Thus the analytic radius of control comes from
$B_K$, while the arithmetic normalization is needed later to turn the
rational evaluation into an integer.  These are separate uses of the
same determinant identity.

\noindent\textbf{[Conditional]} To continue as in source (C.7), assume
the additional input
\begin{equation}
 0<Q_K(\xi)\leq e^{-\varepsilon K^2},\qquad \varepsilon>0,
 \quad\text{for all sufficiently large allowed }K.
 \label{c-negative-hypothesis}
\end{equation}
Choose constants $c>\lambda/\varepsilon$ and $\nu>\rho c$.  For a
sufficiently large positive integer $b$, take
$K=40\lceil c\log b/40\rceil$, so $c\log b\leq K<c\log b+40$.
If $|a/b-\xi|\leq b^{-\nu}$, put $\delta=a/b-\xi$.  The relative norm
bound gives
\[
 |\delta|\|B_K\|\leq e^{40\rho}b^{\rho c-\nu}=o(K^{-1}).
\]
The determinant in \eqref{c-determinant-perturbation} is consequently
positive and at most two eventually, uniformly in $a$.  Integrality
and degree $h$ would then force
\begin{equation}
 \begin{aligned}
 0<b^hQ_K(a/b)
 &\leq2\exp\{\lambda K\log b-\varepsilon K^2\}\\
 &\leq2\exp\{-(\varepsilon-\lambda/c)K^2\}<1,
 \qquad b^hQ_K(a/b)\in\mathbb Z,
 \end{aligned}
 \label{c-rational-contradiction}
\end{equation}
an impossibility.  Under the stated hypotheses this excludes such
approximations for every $\nu>\rho\lambda/\varepsilon$, after choosing
$c$ between $\lambda/\varepsilon$ and $\nu/\rho$.

\noindent\textbf{[Blocked conclusion]} Section~7 supplies
$B_{9,100000}\simeq1.974551870160>0$, not the negative hypothesis
\eqref{c-negative-hypothesis}.  The source's choices
$\varepsilon=79/1600$, $c=75/4$, and $\nu=260$ therefore do not follow
for this ninth-order construction.  The calculation above describes
the consequence of a future negative budget, not an approximation
bound obtained here.

\subsection{Coefficient height and primitive normalization}
\label{c-height-content}

\noindent\textbf{[Derived]} For $Q\in\mathbb R[X]$, let $H(Q)$ be the
largest absolute value of its coefficients.  Since
$1<\xi<2$ and \eqref{c-determinant-perturbation} is a polynomial
identity, the sum of the absolute coefficient values is bounded by
\[
 H(Q_K)\leq Q_K(\xi)
       \prod_{j=1}^{h}\bigl(1+(1+|\xi|)|\beta_j|\bigr)
 \leq Q_K(\xi)(1+3e^{\rho K})^h.
\]
\noindent\textbf{[Conditional]} With the arithmetic and real estimates
already stated, this gives the analogue of source (C.8), without
requiring any additional negative-budget assumption:
\begin{equation}
 \begin{aligned}
 \log H(Q_{9,K,M})
     &\leq(B_{9,M}+\lambda\rho)K^2+o_M(K^2),
       &\lambda\rho&=\frac{256521}{20000}=12.82605,\\
 \log H(Q_K)&\leq(\lambda\rho-\varepsilon)K^2+o(K^2)
       &&\text{if \eqref{c-negative-hypothesis} holds}.
 \end{aligned}
 \label{c-height-bound}
\end{equation}
For $M=100000$ the coefficient in the first line is approximately
$14.800601870160$.  Thus a height estimate remains available, although
the source's height coefficient based on its negative margin is not
transferred.

For nonzero $A\in\mathbb Q[X]$, define $\operatorname{cont}_{\mathbb Q}(A)>0$
so that $A/\operatorname{cont}_{\mathbb Q}(A)$ is an integer polynomial
whose coefficients have greatest common divisor one.  Put
\begin{equation}
 P_K=\frac{\Delta_{9,K}}{\operatorname{cont}_{\mathbb Q}(\Delta_{9,K})},
 \qquad Q_{9,K,M}=a_{K,M}P_K,\qquad a_{K,M}\in\mathbb Z_{>0}.
 \label{c-primitive-normalization}
\end{equation}
The last assertion requires the integrality of $Q_{9,K,M}$ from
Section~5.  To see it, the multiplier is initially a positive rational
number.  A denominator in lowest terms would have to divide every
coefficient of the primitive polynomial $P_K$, so that denominator is
one.  Equivalently, a Bezout combination of its coefficients equals
one and forces the multiplier to be integral.  Hence
\begin{equation}
 0<P_K(\xi)=\frac{Q_{9,K,M}(\xi)}{a_{K,M}}
       \leq Q_{9,K,M}(\xi),\qquad
 H(P_K)=\frac{H(Q_{9,K,M})}{a_{K,M}}\leq H(Q_{9,K,M}).
 \label{c-primitive-bounds}
\end{equation}
The primitive polynomial is independent of the cutoff $M$.  All
available upper bounds descend to it, including any hypothetical
small-value bound.  This removal of content does not itself produce a
negative quadratic coefficient; no size estimate for $a_{K,M}$ has
been obtained that would supply the missing decay.

\clearpage
\begingroup
\small
\section*{Sources and drafting record}
\addcontentsline{toc}{section}{Sources and drafting record}

The structural source for every numbered section and appendix is
Fauzan's manuscript \cite{F}, dated 17 September 2026, version 1,
32 pages. Its equations are cited by their original numbers throughout.
The title below records the title of that manuscript; it is not an
endorsement of its main claim. The downloaded PDF is preserved locally,
with SHA-256
\begin{center}\small\ttfamily
a42c05ccc46b9a7c36428913453c4617202\\
c59109713ba20f464ab9ec640d0da.
\end{center}

\begin{center}
\small
\begin{tabular}{@{}>{\raggedright\arraybackslash}p{0.18\textwidth}>{\raggedright\arraybackslash}p{0.22\textwidth}>{\raggedright\arraybackslash}p{0.50\textwidth}@{}}
\toprule
Draft & Source location & Treatment\\
\midrule
Sections 1--2 & pp.~1--5 & New overview; ninth-order functional, determinant, degree, and positive integral.\\
Section 3 & pp.~5--9 & Ninth-order pullback, completion, distribution, and small-prime calculation.\\
Section 4 & pp.~9--13 & New weights, moment cutoff, and outer costs; inner analytic transfer marked formal.\\
Section 5 & pp.~13--17 & New normalization budget, with dependencies carried forward.\\
Section 6 & pp.~17--20 & Same leading field; new weight bounds and scaling errors.\\
Section 7 & p.~21 & Sign comparison and conditional completion.\\
Appendix A & pp.~21--26 & Same measure and elementary formulas; source's 684 interval checks imported as a certificate.\\
Appendix B & pp.~26--28 & New exact rational costs and revised margins.\\
Appendix C & pp.~28--31 & New scaling and residues; conditional approximation and height calculations.\\
\bottomrule
\end{tabular}
\end{center}

The standard identities used to construct the new functional are
documented in the NIST Digital Library of Mathematical Functions:
Euler's even-zeta formula \cite{D1}, Hermite's integral \cite{D2},
the Bernoulli difference, reflection, and multiplication identities
\cite{D3}, and von Staudt--Clausen \cite{D4}.
The exact rational calculations and finite valuation examples are
calculations for this exploration; they are not external sources.
Earlier local seventh-order notes were used as working context, with
the ninth-order substitutions written out here.

Six separate section assignments used GPT-6 Astra at extra-high
reasoning (\texttt{xhigh}, not \texttt{ultra}): Section~2; Section~3;
Section~4; Section~5 with Appendix~B; Section~6 with Appendix~A; and
Section~7 with Appendix~C. The coordinating agent wrote the introduction,
aligned the shared definitions and constants, and assembled and checked
the document's layout. This division of work is not independent peer
review. Handoff notes and the shared editorial contract are included
with the editable source bundle.

\begin{thebibliography}{9}
\raggedright
\bibitem{F} Aabir Fauzan, \emph{$\zeta(5)$ is irrational},
17 September 2026, version 1, Zenodo.
\href{https://zenodo.org/records/22826419}{doi:10.5281/zenodo.22826419}.
\bibitem{D1} NIST Digital Library of Mathematical Functions,
\href{https://dlmf.nist.gov/25.6.E2}{equation 25.6.2},
even values of the Riemann zeta function.
\bibitem{D2} NIST Digital Library of Mathematical Functions,
\href{https://dlmf.nist.gov/25.11.E29}{equation 25.11.29},
Hermite integral for the Hurwitz zeta function.
\bibitem{D3} NIST Digital Library of Mathematical Functions,
\href{https://dlmf.nist.gov/24.4}{Section 24.4}, particularly
equations 24.4.1, 24.4.3, and 24.4.18.
\bibitem{D4} NIST Digital Library of Mathematical Functions,
\href{https://dlmf.nist.gov/24.10.i}{Section 24.10(i)},
von Staudt--Clausen theorem.
\end{thebibliography}
\noindent\small Source URLs last consulted for this exploration:
14 November 2026. The source bundle includes a machine-readable ledger
of locations, scope, and file hashes.
\endgroup

\end{document}
